\chapter{Radon Measures}\label{radon-measures}
\chapterstatus{complete}

\section{Introduction}\label{radon-measures:section-introduction}

On a locally compact Hausdorff space the measures compatible with the topology are the Radon measures: Borel measures finite on compact sets and approximable from outside by open sets and from inside
by compact sets. They are exactly the positive linear functionals on continuous functions with compact support (the Riesz--Markov theorem), and their signed and vector-valued versions describe the duals
of spaces of continuous functions. This chapter proves these representation theorems, constructs Haar measure on locally compact groups (with the modular function and
Steinhaus' theorem), studies weak convergence of measures and its compactness theorem, and differentiates one Radon measure with
respect to another on $\RR^n$, and with respect to doubling measures on metric spaces; Lebesgue's theorem on the differentiability of monotone functions is a corollary. These results are used in
functional analysis (Chapter~\ref{functional-analysis}), in probability, and in geometric measure theory (Chapter~\ref{geometric-measure-theory}).
References: \cite{rudin-rca}, \cite{folland-real} and \cite{mattila-1995}.

\noindent\textbf{Prerequisites.} Chapter~\ref{general-topology} (General Topology), Chapter~\ref{real-analysis} (Real Analysis), Chapter~\ref{topological-manifolds} (Topological Manifolds and Group Actions),
Chapter~\ref{measure-theory} (Measure and Integration), Chapter~\ref{lebesgue-differentiation} (Differentiation and Absolute Continuity) and Chapter~\ref{signed-measures} (Signed and Complex Measures).

\section{Radon measures and the Riesz representation theorem}\label{radon-measures:section-radon}

A measure assigns numbers to sets; integration against it assigns numbers to functions, $f \mapsto \int f\,d\mu$, and this assignment is linear and positive. On $\RR^n$ the
Riemann integral of continuous functions with compact support is such an assignment, and Lebesgue measure can be recovered from it: $\lambda([0, 1]^n)$ is the infimum of the integrals of the
continuous functions that are at least $1$ on the cube. The Riesz representation theorem says that this always works: every positive linear functional on continuous functions with compact support
is integration against exactly one measure of a reasonable kind. This point of view, \emph{measures as linear functionals}, is basic in geometric measure theory, in probability and in the theory
of partial differential equations. It allows one to take limits of measures (Section~\ref{radon-measures:section-weak-convergence}), and it is the model for distributions, which are linear functionals on
smooth functions (Chapter~\ref{distributions}).

Throughout this section $X$ is a locally compact Hausdorff space with a countable basis; for instance $\RR^n$, an open or a closed subset of $\RR^n$, or a manifold. Such a space is metrisable: it is
regular, because for a point $x$ outside a closed set $F$, Lemma~\ref{general-topology:lemma-lch-shrinking} applied to $\{x\} \subseteq X \setminus F$ gives an open $V \ni x$ with
$\overline{V} \cap F = \emptyset$, and a regular space with a countable basis is metrisable (Theorem~\ref{general-topology:theorem-urysohn-metrisation}). It has an exhaustion by compact sets
$K_1 \subseteq K_2 \subseteq \cdots$ with $K_m \subseteq \operatorname{int}K_{m+1}$ (Proposition~\ref{general-topology:proposition-exhaustion}); every compact $K \subseteq X$ lies in some $K_m$, since the
increasing open sets $\operatorname{int}K_m$ cover $K$.

We write $C_c(X)$ for the real vector space of continuous functions $f \colon X \to \RR$ with compact support, and, for an open $V \subseteq X$,
\[
f \prec V \quad \text{if} \quad f \in C_c(X), \quad 0 \leq f \leq 1 \quad \text{and} \quad \operatorname{supp}f \subseteq V.
\]
Continuous functions are Borel measurable. If $\mu$ is a Borel measure that is finite on compact sets, every $f \in C_c(X)$ is integrable, with $|\int f\,d\mu| \leq \|f\|_\infty\,\mu(\operatorname{supp}f)$.

\begin{lemma}\label{radon-measures:lemma-closed-approximation}
Let $\nu$ be a finite Borel measure on a metrisable space $X$. For every Borel set $E$ and every $\varepsilon > 0$ there are a closed set $F$ and an open set $U$ with $F \subseteq E \subseteq U$ and
$\nu(U \setminus F) < \varepsilon$.
\end{lemma}

\begin{proof}
Fix a metric $d$ inducing the topology, and let $\mathcal{G}$ be the family of Borel sets $E$ with the stated property. A closed set $E$ belongs to $\mathcal{G}$: take $F = E$ and
$U_k = \{x : d(x, E) < 1/k\}$, which are open, decrease, and have intersection $E$ because $E$ is closed; as $\nu$ is finite, $\nu(U_k \setminus E) \to 0$
(Proposition~\ref{measure-theory:proposition-measure-properties}). The family $\mathcal{G}$ is closed under complements, since $F \subseteq E \subseteq U$ gives
$X \setminus U \subseteq X \setminus E \subseteq X \setminus F$ with the same difference $U \setminus F$. It is closed under countable unions: if $F_k \subseteq E_k \subseteq U_k$ with
$\nu(U_k \setminus F_k) < \varepsilon 2^{-k-1}$, then $U = \bigcup_kU_k$ is open and $F^{\prime} = \bigcup_kF_k$ satisfies $\nu(U \setminus F^{\prime}) < \varepsilon/2$; the closed sets
$F_1 \cup \cdots \cup F_N$ increase to $F^{\prime}$, so one of them, $F$, has $\nu(F^{\prime} \setminus F) < \varepsilon/2$, and then $F \subseteq \bigcup_kE_k \subseteq U$ and $\nu(U \setminus F) < \varepsilon$.
So $\mathcal{G}$ is a $\sigma$-algebra containing the closed sets, hence all Borel sets.
\end{proof}

\begin{proposition}\label{radon-measures:proposition-radon-automatic}
Let $X$ be a locally compact Hausdorff space with a countable basis and $\mu$ a Borel measure on $X$ that is finite on compact sets. Then $\mu$ is $\sigma$-finite and regular, hence a Radon measure
(Definition~\ref{measure-theory:definition-regular}): for every Borel set $E$,
\[
\mu(E) = \inf\{\mu(U) : U \supseteq E \text{ open}\} = \sup\{\mu(K) : K \subseteq E \text{ compact}\}.
\]
\end{proposition}

\begin{proof}
$X$ is the union of the compact sets $K_m$ of an exhaustion, which have finite measure, so $\mu$ is $\sigma$-finite; and every point has a neighbourhood $\operatorname{int}K_{m+1} \subseteq K_{m+1}$ of finite
measure. The measures $\nu_m(A) = \mu(A \cap \operatorname{int}K_m)$ are finite, and Lemma~\ref{radon-measures:lemma-closed-approximation} applies to them, $X$ being metrisable.

\emph{Outer regularity.} Let $E$ be a Borel set with $\mu(E) < \infty$ (otherwise there is nothing to prove) and $\varepsilon > 0$. Choose open sets $U_m \supseteq E$ with
$\nu_m(U_m \setminus E) < \varepsilon2^{-m}$, and let $U = \bigcup_m(U_m \cap \operatorname{int}K_m)$. It is open and contains $E = \bigcup_m(E \cap \operatorname{int}K_m)$, and
$\mu(U \setminus E) \leq \sum_m\mu\bigl((U_m \cap \operatorname{int}K_m) \setminus E\bigr) = \sum_m\nu_m(U_m \setminus E) < \varepsilon$. So $\mu(U) < \mu(E) + \varepsilon$.

\emph{Inner regularity.} Let $E$ be a Borel set and $t < \mu(E)$. By continuity from below, $\mu(E \cap \operatorname{int}K_m) > t$ for some $m$. Choose a closed $F \subseteq E$ with
$\nu_m(E \setminus F) < \mu(E \cap \operatorname{int}K_m) - t$. Then $K = F \cap K_m$ is compact, contained in $E$, and
$\mu(K) \geq \mu(F \cap \operatorname{int}K_m) = \nu_m(E) - \nu_m(E \setminus F) > t$.
\end{proof}

So on $\RR^n$ and on manifolds, \emph{Radon measure} simply means \emph{Borel measure finite on compact sets}. For general locally compact Hausdorff spaces this is false, and the theory below has to
be set up with more care \cite[Chapter 2]{rudin-rca}; we shall not need that generality.

\begin{lemma}\label{radon-measures:lemma-partition-compact}
Let $X$ be a locally compact Hausdorff space, $K \subseteq X$ compact and $V_1, \ldots, V_n$ open sets covering $K$. There are $h_i \prec V_i$ with $\sum_ih_i \leq 1$ on $X$ and $\sum_ih_i = 1$ on $K$.
\end{lemma}

\begin{proof}
Each $x \in K$ lies in some $V_i$, and Lemma~\ref{general-topology:lemma-lch-shrinking} gives an open $W_x \ni x$ with compact closure $\overline{W_x} \subseteq V_i$. Finitely many $W_x$ cover $K$; let $H_i$ be
the union of the closures of those among them for which the chosen $i$ was $i$. Then $H_i \subseteq V_i$ is compact and $K \subseteq H_1 \cup \cdots \cup H_n$. Lemma~\ref{general-topology:lemma-lch-shrinking}
gives $g_i \prec V_i$ with $g_i = 1$ on $H_i$. Put $h_1 = g_1$ and $h_k = (1 - g_1)\cdots(1 - g_{k-1})g_k$. Then $0 \leq h_k \leq g_k$, so $h_k \prec V_k$, and by induction on $m$,
\[
h_1 + \cdots + h_m = 1 - (1 - g_1)(1 - g_2)\cdots(1 - g_m),
\]
since $1 - \prod_{k \leq m}(1 - g_k) + h_{m+1} = 1 - \prod_{k \leq m+1}(1 - g_k)$. For $m = n$ the right side lies in $[0, 1]$, and it equals $1$ on each $H_i$, where the factor $1 - g_i$ vanishes.
\end{proof}

\begin{definition}\label{radon-measures:definition-positive-functional}
A linear map $\Lambda \colon C_c(X) \to \RR$ is \emph{positive} if $\Lambda(f) \geq 0$ whenever $f \geq 0$.
\end{definition}

A positive functional is monotone: $f \leq g$ implies $\Lambda f \leq \Lambda g$. It is also automatically continuous in the following sense. If $K$ is compact and $\phi_K \in C_c(X)$ satisfies
$0 \leq \phi_K \leq 1$ and $\phi_K = 1$ on $K$ (Lemma~\ref{general-topology:lemma-lch-shrinking}), then every $f$ with $\operatorname{supp}f \subseteq K$ satisfies
$-\|f\|_\infty\phi_K \leq f \leq \|f\|_\infty\phi_K$, hence
\[
|\Lambda f| \leq \Lambda(\phi_K)\,\|f\|_\infty .
\]
So uniform convergence of functions supported in a fixed compact set implies convergence of the values of $\Lambda$.

\begin{example}\label{radon-measures:example-positive-functionals}
\begin{enumerate}
\item \emph{Point masses.} $\Lambda f = f(x_0)$ is positive, and $\Lambda f = \int f\,d\delta_{x_0}$. More generally, for points $x_1, x_2, \ldots$ of $X$ with no accumulation point and numbers
$c_k \geq 0$, the finite sum $\Lambda f = \sum_kc_kf(x_k)$ is integration against $\sum_kc_k\delta_{x_k}$.
\item \emph{Densities.} For $g \in L^1_{\mathrm{loc}}(\RR^n)$ with $g \geq 0$, $\Lambda f = \int fg\,d\lambda$ is positive; it is integration against the measure $g\lambda$, $E \mapsto \int_Eg\,d\lambda$. For
$g = 1$ this is the Riemann integral of $f$ (Proposition~\ref{measure-theory:proposition-riemann-lebesgue}, together with Fubini's theorem for the iterated integrals).
\item \emph{Arc length on a circle.} On $X = \RR^2$, $\Lambda f = \int_0^{2\pi}f(\cos t, \sin t)\,dt$ is positive. The corresponding measure $\sigma$ (Theorem~\ref{radon-measures:theorem-riesz-markov}
below) lives on the unit circle $S^1$: if $K$ is compact and disjoint from $S^1$, a function $f \prec \RR^2 \setminus S^1$ with $f = 1$ on $K$ gives $\sigma(K) \leq \int f\,d\sigma = \Lambda f = 0$, and
by inner regularity $\sigma(\RR^2 \setminus S^1) = 0$. Since $\lambda^2(S^1) = 0$, the measures $\sigma$ and $\lambda^2$ are mutually singular. The measure of an arc is its length
(Exercise~\ref{radon-measures:exercise-riesz-circle}), and $\sigma$ is the restriction to $S^1$ of the one-dimensional Hausdorff measure of Chapter~\ref{geometric-measure-theory}.
\item \emph{Push-forward.} More generally, for a continuous map $\varphi \colon [a, b] \to X$, the functional $\Lambda f = \int_a^bf(\varphi(t))\,dt$ is integration against the push-forward
$\varphi_*\lambda$, $E \mapsto \lambda(\{t \in [a, b] : \varphi(t) \in E\})$, which is a Borel measure, finite because $\lambda([a, b]) < \infty$.
\end{enumerate}
\end{example}

\begin{theorem}[Riesz--Markov]\label{radon-measures:theorem-riesz-markov}
Let $X$ be a locally compact Hausdorff space with a countable basis, and $\Lambda \colon C_c(X) \to \RR$ a positive linear functional. There is a unique Borel measure $\mu$ on $X$, finite on compact
sets, with
\[
\Lambda f = \int_Xf\,d\mu \qquad \text{for all } f \in C_c(X).
\]
It is a Radon measure, and for open $V$ and compact $K$,
\[
\mu(V) = \sup\{\Lambda f : f \prec V\}, \qquad \mu(K) = \inf\{\Lambda f : f \in C_c(X),\ f \geq \mathbf{1}_K\}.
\]
\end{theorem}

\begin{proof}
The proof follows \cite[Theorem 2.14]{rudin-rca}. Motivated by the formulas to be proved, define $\mu(V) = \sup\{\Lambda f : f \prec V\}$ for open $V$, and for an arbitrary $E \subseteq X$
\[
\mu^*(E) = \inf\{\mu(V) : V \supseteq E \text{ open}\}.
\]
If $V \subseteq W$ are open, $f \prec V$ implies $f \prec W$, so $\mu(V) \leq \mu(W)$; hence $\mu^*(V) = \mu(V)$ for open $V$, and $\mu^*$ is monotone.

\emph{Step 1: $\mu^*$ is an outer measure.} First let $V = \bigcup_iV_i$ be a countable union of open sets, and $f \prec V$. The compact set $\operatorname{supp}f$ is covered by finitely many
$V_1, \ldots, V_n$, and Lemma~\ref{radon-measures:lemma-partition-compact} gives $h_i \prec V_i$ with $\sum_ih_i = 1$ on $\operatorname{supp}f$. Then $f = \sum_ih_if$ with $h_if \prec V_i$, so
$\Lambda f = \sum_i\Lambda(h_if) \leq \sum_i\mu(V_i)$; taking the supremum over $f$, $\mu(V) \leq \sum_i\mu(V_i)$. Now let $E = \bigcup_iE_i$ be arbitrary, with $\sum_i\mu^*(E_i) < \infty$, and
$\varepsilon > 0$. Choose open $V_i \supseteq E_i$ with $\mu(V_i) \leq \mu^*(E_i) + \varepsilon2^{-i}$; then
$\mu^*(E) \leq \mu(\bigcup_iV_i) \leq \sum_i\mu^*(E_i) + \varepsilon$. Finally $\mu^*(\emptyset) = 0$, as $f \prec \emptyset$ forces $f = 0$.

\emph{Step 2: compact sets.} Let $K$ be compact. If $f \in C_c(X)$ and $f \geq \mathbf{1}_K$, and $0 < \alpha < 1$, then $V_\alpha = \{f > \alpha\}$ is an open set containing $K$, and every $g \prec V_\alpha$
satisfies $\alpha g \leq f$; so $\mu^*(K) \leq \mu(V_\alpha) \leq \Lambda f/\alpha$, and letting $\alpha \to 1$, $\mu^*(K) \leq \Lambda f$. Conversely, for an open $V \supseteq K$,
Lemma~\ref{general-topology:lemma-lch-shrinking} gives $f \prec V$ with $f = 1$ on $K$, and $\Lambda f \leq \mu(V)$. So $\mu^*(K) = \inf\{\Lambda f : f \geq \mathbf{1}_K\}$, which is finite.

\emph{Step 3: Borel sets are $\mu^*$-measurable.} By Carath\'eodory's theorem (Theorem~\ref{measure-theory:theorem-caratheodory}) the $\mu^*$-measurable sets form a $\sigma$-algebra on which $\mu^*$ is a
measure, so it suffices to show that every open $U$ is $\mu^*$-measurable: $\mu^*(E \cap U) + \mu^*(E \setminus U) \leq \mu^*(E)$ for every $E$ with $\mu^*(E) < \infty$. Let $\varepsilon > 0$ and
choose an open $V \supseteq E$ with $\mu(V) \leq \mu^*(E) + \varepsilon$. Choose $f \prec V \cap U$ with $\Lambda f > \mu(V \cap U) - \varepsilon$; its support $K$ is a compact subset of $U$, so
$V \setminus K$ is open and we may choose $g \prec V \setminus K$ with $\Lambda g > \mu(V \setminus K) - \varepsilon$. The functions $f$ and $g$ have disjoint supports, so $f + g \prec V$, and
\[\begin{aligned}
\mu^*(E) + \varepsilon \geq \mu(V) \geq \Lambda f + \Lambda g &> \mu(V \cap U) + \mu(V \setminus K) - 2\varepsilon \\
&\geq \mu^*(E \cap U) + \mu^*(E \setminus U) - 2\varepsilon,
\end{aligned}\]
because $E \cap U \subseteq V \cap U$ and $E \setminus U \subseteq V \setminus K$. Let $\varepsilon \to 0$.

So the restriction of $\mu^*$ to the Borel sets is a Borel measure $\mu$, finite on compact sets by Step 2, and by construction it satisfies the two formulas of the statement.

\emph{Step 4: $\Lambda f = \int f\,d\mu$.} Both sides are linear in $f$, so it suffices to show $\Lambda f \leq \int f\,d\mu$ for every $f \in C_c(X)$: applied to $-f$ it gives the reverse inequality.
Let $K = \operatorname{supp}f$, choose $a < b$ with $f(X) \subseteq [a, b]$, let $\varepsilon > 0$, and choose $y_0 < a < y_1 < \cdots < y_n = b$ with $y_i - y_{i-1} < \varepsilon$. The Borel sets
$E_i = \{y_{i-1} < f \leq y_i\} \cap K$ are disjoint with union $K$. By the definition of $\mu^*$ there are open sets $V_i \supseteq E_i$ with $\mu(V_i) < \mu(E_i) + \varepsilon/n$; replacing $V_i$ by
$V_i \cap \{f < y_i + \varepsilon\}$ we may assume $f < y_i + \varepsilon$ on $V_i$. Lemma~\ref{radon-measures:lemma-partition-compact} gives $h_i \prec V_i$ with $\sum_ih_i = 1$ on $K$. Then
$f = \sum_ih_if$, and $h_if \leq (y_i + \varepsilon)h_i$. By Step 2, $\sum_i\Lambda h_i = \Lambda(\sum_ih_i) \geq \mu(K)$, and $\Lambda h_i \leq \mu(V_i)$. As $|a| + y_i + \varepsilon > 0$,
\[
\begin{aligned}
\Lambda f = \sum_i\Lambda(h_if) &\leq \sum_i(y_i + \varepsilon)\Lambda h_i = \sum_i(|a| + y_i + \varepsilon)\Lambda h_i - |a|\sum_i\Lambda h_i\\
&\leq \sum_i(|a| + y_i + \varepsilon)\Bigl(\mu(E_i) + \frac{\varepsilon}{n}\Bigr) - |a|\mu(K)\\
&\leq \sum_i(y_i - \varepsilon)\mu(E_i) + 2\varepsilon\mu(K) + \varepsilon(|a| + b + \varepsilon),
\end{aligned}
\]
using $\sum_i\mu(E_i) = \mu(K)$ and $y_i \leq b$. On $E_i$ we have $f > y_{i-1} > y_i - \varepsilon$, so $\sum_i(y_i - \varepsilon)\mu(E_i) \leq \int f\,d\mu$. Letting $\varepsilon \to 0$ gives
$\Lambda f \leq \int f\,d\mu$.

\emph{Step 5: regularity and uniqueness.} $\mu$ is a Radon measure by Proposition~\ref{radon-measures:proposition-radon-automatic}. Let $\mu^{\prime}$ be another Borel measure, finite on compact sets,
representing $\Lambda$, and $K$ compact. If $f \geq \mathbf{1}_K$, then $\mu^{\prime}(K) \leq \int f\,d\mu^{\prime} = \Lambda f$, so $\mu^{\prime}(K) \leq \mu(K)$ by Step 2. If $V \supseteq K$ is open and
$f \prec V$ with $f = 1$ on $K$, then $\mu(K) \leq \Lambda f = \int f\,d\mu^{\prime} \leq \mu^{\prime}(V)$, and by outer regularity of $\mu^{\prime}$
(Proposition~\ref{radon-measures:proposition-radon-automatic}) $\mu(K) \leq \mu^{\prime}(K)$. So $\mu$ and $\mu^{\prime}$ agree on compact sets, hence, by inner regularity, on all Borel sets.
\end{proof}

\begin{remark}\label{radon-measures:remark-riesz-markov}
\begin{enumerate}
\item For $X = \RR^n$ and $\Lambda$ the Riemann integral, the theorem produces Lebesgue measure on the Borel sets; the uniqueness statement is a convenient way to identify measures, as in
Exercise~\ref{radon-measures:exercise-riesz-circle}. On a smooth manifold, an orientation and a volume form, or more generally a density, define a positive functional by integration, and hence a
Radon measure (Chapter~\ref{differential-forms}).
\item The theorem holds for every locally compact Hausdorff space, but then the measure is only outer regular and inner regular on open sets and on sets of finite measure
\cite[Theorem 2.14 and Theorem 2.18]{rudin-rca}.
\item If $X$ is compact, $1 \in C_c(X)$, and $|\Lambda f| \leq \Lambda(1)\|f\|_\infty = \mu(X)\|f\|_\infty$: a positive functional on $C(X)$ is bounded with norm $\mu(X)$. Every bounded linear functional on
$C(X)$ is integration against a \emph{signed} Radon measure, a difference of two finite Radon measures \cite[Theorem 6.19]{rudin-rca} (for the spaces of this section, see Corollary~\ref{radon-measures:corollary-signed-riesz}
below); this identifies the dual space of $C(X)$, a theme of
Chapter~\ref{functional-analysis}.
\end{enumerate}
\end{remark}

\subsection{Signed and vector-valued functionals}\label{radon-measures:subsection-riesz-vector}

Many natural linear functionals are not positive. The functional $\varphi \mapsto \varphi(0) - \varphi(1)$ on $C_c(\RR)$ is the difference of two Dirac measures; $\varphi \mapsto \int\varphi h\,d\lambda$ with $h$
changing sign is the difference of the measures $h^+\lambda$ and $h^-\lambda$. In geometric measure theory one also meets functionals on \emph{vector fields}: for the unit circle $S^1 \subseteq \RR^2$ with
outward normal $n(x) = x$, the \emph{flux} $F \mapsto \int_0^{2\pi}\langle F(\cos t, \sin t), (\cos t, \sin t)\rangle\,dt$ through the circle is a linear functional on continuous vector fields $F \colon \RR^2 \to \RR^2$ with compact support. The
Riesz representation theorem extends to all these cases: a functional that is merely bounded on functions of bounded size supported in a fixed compact set is integration against a measure, after pairing with
a field of unit vectors.

Let $H$ be a finite-dimensional real vector space with an inner product $\langle\,,\rangle$ and norm $|\ |$; for instance $H = \RR^m$ with the standard inner product. Let $C_c(X, H)$ be the vector space of
continuous maps $f \colon X \to H$ with compact support. For $f \in C_c(X, H)$ the function $|f|$ belongs to $C_c(X)$. The map $r \colon H \to H$, $r(v) = v$ for $|v| \leq 1$ and $r(v) = v/|v|$ for $|v| > 1$, the
nearest point projection onto the closed unit ball, is continuous, takes values in the ball, and moves no point farther from the ball: $|r(u) - w| \leq |u - w|$ whenever $|w| \leq 1$. To see this,
we may assume $|u| > 1$; then
$|u - w|^2 - |r(u) - w|^2 = |u|^2 - 1 - 2(|u| - 1)\langle u/|u|, w\rangle \geq |u|^2 - 1 - 2(|u| - 1) = (|u| - 1)^2 \geq 0$.

\begin{lemma}\label{radon-measures:lemma-radon-approximation}
Let $\mu$ be a Radon measure on $X$.
\begin{enumerate}
\item For $1 \leq p < \infty$, $C_c(X)$ is dense in $L^p(\mu)$.
\item Let $w \colon X \to H$ be Borel with $|w| \leq 1$, $E$ a Borel set with $\mu(E) < \infty$, and $V \supseteq E$ open. For every $\varepsilon > 0$ there is $g \in C_c(X, H)$ with $|g| \leq 1$,
$\operatorname{supp}g \subseteq V$ and $\int|g - \mathbf{1}_Ew|\,d\mu < \varepsilon$.
\end{enumerate}
\end{lemma}

\begin{proof}
(1) As in Proposition~\ref{measure-theory:proposition-density}, simple functions $\sum a_i\mathbf{1}_{E_i}$ with $\mu(E_i) < \infty$ are dense in $L^p(\mu)$, by
Proposition~\ref{measure-theory:proposition-measurable-functions}(4) and dominated convergence. Given $E$ with $\mu(E) < \infty$ and $\varepsilon > 0$, regularity of $\mu$ gives a compact $K$ and an
open $U$ with $K \subseteq E \subseteq U$ and $\mu(U \setminus K) < \varepsilon$, and Lemma~\ref{general-topology:lemma-lch-shrinking} gives $\phi \in C_c(X)$ with $0 \leq \phi \leq 1$, $\phi = 1$ on $K$ and
$\operatorname{supp}\phi \subseteq U$. Then $\int|\mathbf{1}_E - \phi|^p\,d\mu \leq \mu(U \setminus K) < \varepsilon$.

(2) Fix an orthonormal basis $e_1, \ldots, e_m$ of $H$ and write $w = \sum_jw_je_j$. By (1) choose $\phi_j \in C_c(X)$ with $\int|\phi_j - \mathbf{1}_Ew_j|\,d\mu < \varepsilon/(2m)$, and put
$g_0 = \sum_j\phi_je_j$, so $\int|g_0 - \mathbf{1}_Ew|\,d\mu < \varepsilon/2$. Then $g_1 = r \circ g_0$ is continuous with compact support, $|g_1| \leq 1$, and $|g_1 - \mathbf{1}_Ew| \leq |g_0 - \mathbf{1}_Ew|$
pointwise, as $|\mathbf{1}_Ew| \leq 1$. Choose a compact $K \subseteq E$ with $\mu(E \setminus K) < \varepsilon/4$ and $\psi \in C_c(X)$ with $0 \leq \psi \leq 1$, $\psi = 1$ on $K$ and
$\operatorname{supp}\psi \subseteq V$, and let $g = \psi g_1$. On $K$, $g = g_1$; outside $E$, $|g - \mathbf{1}_Ew| = |g| \leq |g_1| = |g_1 - \mathbf{1}_Ew|$; on $E \setminus K$, $|g - w| \leq 2$. Hence
$\int|g - \mathbf{1}_Ew|\,d\mu \leq \int|g_1 - \mathbf{1}_Ew|\,d\mu + 2\mu(E \setminus K) < \varepsilon$.
\end{proof}

\begin{theorem}[Riesz representation theorem for vector-valued functions]\label{radon-measures:theorem-riesz-vector}
Let $L \colon C_c(X, H) \to \RR$ be linear and suppose that for every compact $K \subseteq X$
\[
M_K = \sup\bigl\{L(f) : f \in C_c(X, H),\ |f| \leq 1,\ \operatorname{supp}f \subseteq K\bigr\} < \infty .
\]
Then there are a Radon measure $\mu$ on $X$ and a Borel map $\nu \colon X \to H$ with $|\nu| = 1$ $\mu$-almost everywhere such that
\[
L(f) = \int_X\langle f, \nu\rangle\,d\mu \qquad\text{for all } f \in C_c(X, H).
\]
Moreover, for every open $V \subseteq X$,
\begin{equation}\label{radon-measures:equation-total-variation}
\mu(V) = \sup\bigl\{L(f) : f \in C_c(X, H),\ |f| \leq 1,\ \operatorname{supp}f \subseteq V\bigr\},
\end{equation}
and the pair $(\mu, \nu)$ is unique: $\mu$ is unique, and $\nu$ is unique up to $\mu$-null sets. The measure $\mu$ is the \emph{total variation measure} of $L$.
\end{theorem}

\begin{proof}
The idea is to manufacture a positive functional out of $L$, namely the largest value of $L$ on fields dominated by a given function, and to apply the Riesz--Markov theorem to it.

\emph{Step 1: a positive functional on nonnegative functions.} For $\varphi \in C_c(X)$ with $\varphi \geq 0$ let
\[
\lambda(\varphi) = \sup\bigl\{L(g) : g \in C_c(X, H),\ |g| \leq \varphi\bigr\}.
\]
Taking $g = 0$ shows $\lambda(\varphi) \geq 0$. If $|g| \leq \varphi$, then $\operatorname{supp}g \subseteq K = \operatorname{supp}\varphi$, and if $\varphi \neq 0$, $g/\|\varphi\|_\infty$ is admissible in the
definition of $M_K$; so $L(g) \leq \|\varphi\|_\infty M_K$ and $\lambda(\varphi) < \infty$. Clearly $\lambda(c\varphi) = c\lambda(\varphi)$ for $c > 0$, as $g \mapsto cg$ matches the admissible fields. We
show $\lambda(\varphi_1 + \varphi_2) = \lambda(\varphi_1) + \lambda(\varphi_2)$ for $\varphi_1, \varphi_2 \geq 0$. If $|g_i| \leq \varphi_i$, then $|g_1 + g_2| \leq \varphi_1 + \varphi_2$, so
$L(g_1) + L(g_2) \leq \lambda(\varphi_1 + \varphi_2)$; taking suprema gives ``$\geq$''. Conversely let $|g| \leq \varphi_1 + \varphi_2$, and define $g_i = \varphi_ig/(\varphi_1 + \varphi_2)$ where
$\varphi_1 + \varphi_2 > 0$ and $g_i = 0$ elsewhere. Then $|g_i| \leq \varphi_i$, and $g_i$ is continuous: this is clear where $\varphi_1 + \varphi_2 > 0$, and at a point $x_0$ where
$\varphi_1 + \varphi_2 = 0$ we have $|g_i| \leq \varphi_i$ everywhere and $\varphi_i(x_0) = 0$, so $g_i(x) \to 0 = g_i(x_0)$ as $x \to x_0$. Since $g = 0$ where $\varphi_1 + \varphi_2 = 0$, $g = g_1 + g_2$, and
$L(g) = L(g_1) + L(g_2) \leq \lambda(\varphi_1) + \lambda(\varphi_2)$.

\emph{Step 2: extension and the measure $\mu$.} For $\varphi \in C_c(X)$ put $\lambda(\varphi) = \lambda(\varphi^+) - \lambda(\varphi^-)$, where $\varphi^\pm = \max(\pm\varphi, 0) \in C_c(X)$. If
$\varphi = a - b$ with $a, b \geq 0$ in $C_c(X)$, then $\varphi^+ + b = \varphi^- + a$, and Step 1 gives $\lambda(\varphi) = \lambda(a) - \lambda(b)$. Writing
$\varphi + \psi = (\varphi^+ + \psi^+) - (\varphi^- + \psi^-)$ shows that $\lambda$ is additive; it is homogeneous for $c \geq 0$ by Step 1, and $\lambda(-\varphi) = -\lambda(\varphi)$ since
$(-\varphi)^\pm = \varphi^\mp$. So $\lambda$ is a positive linear functional on $C_c(X)$, and the Riesz--Markov theorem (Theorem~\ref{radon-measures:theorem-riesz-markov}) gives a Radon measure $\mu$ with
$\lambda(\varphi) = \int\varphi\,d\mu$.

\emph{Step 3: $|L(g)| \leq \int|g|\,d\mu$.} For $g \in C_c(X, H)$, $g$ is admissible for $\lambda(|g|)$, so $L(g) \leq \lambda(|g|) = \int|g|\,d\mu$; apply this also to $-g$.

\emph{Step 4: construction of $\nu$.} Fix an orthonormal basis $e_1, \ldots, e_m$ of $H$ and let $L_j(\varphi) = L(\varphi e_j)$ for $\varphi \in C_c(X)$. By Step 3, $|L_j(\varphi)| \leq \int|\varphi|\,d\mu$,
so $\Lambda_j(\varphi) = \int\varphi\,d\mu + L_j(\varphi)$ is a positive linear functional on $C_c(X)$. By the Riesz--Markov theorem there is a Radon measure $\sigma_j$ with $\Lambda_j(\varphi) = \int\varphi\,d\sigma_j$.
For open $V$ and $\varphi \prec V$, $\Lambda_j(\varphi) \leq 2\int\varphi\,d\mu \leq 2\mu(V)$, so $\sigma_j(V) \leq 2\mu(V)$ by the formula for $\sigma_j(V)$ in Theorem~\ref{radon-measures:theorem-riesz-markov}; by
outer regularity of $\mu$, $\sigma_j(E) \leq 2\mu(E)$ for all Borel $E$. Both measures are $\sigma$-finite (Proposition~\ref{radon-measures:proposition-radon-automatic}) and $\sigma_j \ll \mu$, so by the
Radon--Nikodym theorem (Theorem~\ref{lebesgue-differentiation:theorem-radon-nikodym}) $\sigma_j(E) = \int_Eh_j\,d\mu$ for a Borel function $h_j \geq 0$. Since $\int_Eh_j\,d\mu \leq 2\mu(E)$ for all $E$, applied to
$E = \{h_j > 2\} \cap K_m$ for the compact exhaustion $(K_m)$, $h_j \leq 2$ almost everywhere, and we may change $h_j$ on a null set so that $0 \leq h_j \leq 2$ everywhere. By the usual extension from
indicator functions to simple functions and limits, $\int\varphi\,d\sigma_j = \int\varphi h_j\,d\mu$ for $\varphi \in C_c(X)$. Put $\nu_j = h_j - 1 \in [-1, 1]$ and $\nu = \sum_j\nu_je_j$, a Borel map with
$|\nu| \leq \sqrt m$. Then $L_j(\varphi) = \int\varphi\,\nu_j\,d\mu$, and for $f \in C_c(X, H)$, writing $f = \sum_j\langle f, e_j\rangle e_j$ with $\langle f, e_j\rangle \in C_c(X)$,
\[
L(f) = \sum_jL_j\bigl(\langle f, e_j\rangle\bigr) = \int\sum_j\langle f, e_j\rangle\nu_j\,d\mu = \int\langle f, \nu\rangle\,d\mu .
\]

\emph{Step 5: $|\nu| \leq 1$ almost everywhere.} Let $E = \{|\nu| > 1\} \cap K_m$ and $w = \nu/|\nu|$ on $E$, $w = 0$ elsewhere. For $\varepsilon > 0$, Lemma~\ref{radon-measures:lemma-radon-approximation}(2) gives
$g \in C_c(X, H)$ with $|g| \leq 1$ and $\int|g - \mathbf{1}_Ew|\,d\mu < \varepsilon$. Then, by Step 3,
\[
\int_E|\nu|\,d\mu = \int\langle\mathbf{1}_Ew, \nu\rangle\,d\mu \leq L(g) + \sqrt m\,\varepsilon \leq \int|g|\,d\mu + \sqrt m\,\varepsilon \leq \mu(E) + (1 + \sqrt m)\varepsilon .
\]
So $\int_E(|\nu| - 1)\,d\mu \leq 0$, while $|\nu| - 1 > 0$ on $E$; hence $\mu(E) = 0$, for every $m$.

\emph{Step 6: $|\nu| \geq 1$ almost everywhere.} For $\varphi \geq 0$ in $C_c(X)$ and $|g| \leq \varphi$, $L(g) = \int\langle g, \nu\rangle\,d\mu \leq \int\varphi|\nu|\,d\mu$; taking the supremum,
$\int\varphi\,d\mu = \lambda(\varphi) \leq \int\varphi|\nu|\,d\mu$, that is, $\int\varphi(1 - |\nu|)\,d\mu \leq 0$. Given a compact $K$, choose such $\varphi$ with $\varphi = 1$ on $K$ and $0 \leq \varphi \leq 1$
(Lemma~\ref{general-topology:lemma-lch-shrinking}); as $1 - |\nu| \geq 0$ almost everywhere by Step 5, $\int_K(1 - |\nu|)\,d\mu \leq \int\varphi(1 - |\nu|)\,d\mu \leq 0$, so $|\nu| = 1$ almost
everywhere on $K$, and on $X = \bigcup_mK_m$.

\emph{Step 7: the formula \eqref{radon-measures:equation-total-variation} and uniqueness.} Let $(\mu^{\prime}, \nu^{\prime})$ be any pair with the properties of the theorem; we show that $\mu^{\prime}(V)$ equals
the right side $s(V)$ of \eqref{radon-measures:equation-total-variation}. If $|f| \leq 1$, then $L(f) = \int\langle f, \nu^{\prime}\rangle\,d\mu^{\prime} \leq \int_V|f|\,d\mu^{\prime} \leq \mu^{\prime}(V)$ for
$\operatorname{supp}f \subseteq V$, so $s(V) \leq \mu^{\prime}(V)$. Conversely, let $K \subseteq V$ be compact and $\varepsilon > 0$. Lemma~\ref{radon-measures:lemma-radon-approximation}(2), applied to $\mu^{\prime}$,
$w = \nu^{\prime}$ and $E = K$, gives $g$ with $|g| \leq 1$, $\operatorname{supp}g \subseteq V$ and $\int|g - \mathbf{1}_K\nu^{\prime}|\,d\mu^{\prime} < \varepsilon$, so
$L(g) \geq \int_K|\nu^{\prime}|^2\,d\mu^{\prime} - \varepsilon = \mu^{\prime}(K) - \varepsilon$. By inner regularity $s(V) \geq \mu^{\prime}(V)$. Hence $\mu^{\prime}(V) = s(V)$ for all open $V$; this proves
\eqref{radon-measures:equation-total-variation} for $\mu$, and shows that $\mu^{\prime}$ and $\mu$ agree on open sets, hence on all Borel sets by outer regularity. Finally
$\int\langle f, \nu - \nu^{\prime}\rangle\,d\mu = 0$ for all $f \in C_c(X, H)$. Let $E = \{\nu \neq \nu^{\prime}\} \cap K_m$ and $w = (\nu - \nu^{\prime})/|\nu - \nu^{\prime}|$ on $E$, $0$ elsewhere; approximating
$\mathbf{1}_Ew$ by $g$ as in Lemma~\ref{radon-measures:lemma-radon-approximation}(2) and using $|\nu - \nu^{\prime}| \leq 2$ almost everywhere gives $\int_E|\nu - \nu^{\prime}|\,d\mu \leq 2\varepsilon$ for every
$\varepsilon$, so $\nu = \nu^{\prime}$ almost everywhere on $E$, and on $X$.
\end{proof}

For a positive functional on $C_c(X) = C_c(X, \RR)$ the theorem gives back the Riesz--Markov measure with $\nu = 1$. For general functionals on $C_c(X)$ it gives the decomposition into a positive and a negative
part.

\begin{corollary}\label{radon-measures:corollary-signed-riesz}
Let $L \colon C_c(X) \to \RR$ be linear with $\sup\{|L(\varphi)| : |\varphi| \leq 1, \operatorname{supp}\varphi \subseteq K\} < \infty$ for every compact $K$. Then there are Radon measures $\mu^+$ and
$\mu^-$, concentrated on disjoint Borel sets $X^+$ and $X^-$, with
\[
L(\varphi) = \int\varphi\,d\mu^+ - \int\varphi\,d\mu^- \qquad (\varphi \in C_c(X)).
\]
If $X$ is compact, every bounded linear functional $L$ on $C(X)$ (with the norm $\|\ \|_\infty$) has this form with $\mu^\pm$ finite, and
\[
\|L\| = \|\mu^+ - \mu^-\| = \mu^+(X) + \mu^-(X),
\]
the \emph{total variation norm} of the signed measure $\mu^+ - \mu^-$, $E \mapsto \mu^+(E) - \mu^-(E)$: the norm of the functional is the norm of the measure.
\end{corollary}

\begin{proof}
Apply Theorem~\ref{radon-measures:theorem-riesz-vector} with $H = \RR$. Then $\nu = \pm1$ almost everywhere; let $X^\pm = \{\nu = \pm1\}$ and $\mu^\pm(E) = \mu(E \cap X^\pm)$, Borel measures finite on compact
sets, hence Radon by Proposition~\ref{radon-measures:proposition-radon-automatic}. Then $\int\varphi\nu\,d\mu = \int\varphi\,d\mu^+ - \int\varphi\,d\mu^-$. If $X$ is compact and $L$ is bounded, the hypothesis
holds with the constant $\|L\|$, $\mu^\pm(X) \leq \mu(X) < \infty$, and \eqref{radon-measures:equation-total-variation} with $V = X$ gives $\mu(X) = \sup\{L(\varphi) : \|\varphi\|_\infty \leq 1\} = \|L\|$.
\end{proof}

For a signed measure $\mu = \mu^+ - \mu^-$ as in the corollary, with $\|\mu\| = \mu^+(X) + \mu^-(X)$, the total variation norm is intrinsic to the set function $\mu$:
$\|\mu\| = \sup\sum_i|\mu(E_i)|$ over all finite partitions of $X$ into Borel sets $E_i$. Indeed $\sum_i|\mu(E_i)| \leq \sum_i(\mu^+(E_i) + \mu^-(E_i)) = \|\mu\|$, with equality for the partition
$\{X^+, X \setminus X^+\}$. For example $\|\delta_0 - \delta_1\| = 2$, and $\|h\lambda\| = \int|h|\,d\lambda$ for the measure $E \mapsto \int_Eh\,d\lambda$ with $h \in L^1$.
So the dual space of $C(X)$, for a compact metrisable $X$, is the space of signed Radon measures with the total variation norm; this proves the statement quoted in
Remark~\ref{radon-measures:remark-riesz-markov}(3) for the spaces of this section. The decomposition $X = X^+ \cup X^-$ is a Hahn decomposition
(Theorem~\ref{signed-measures:theorem-hahn}) and $\mu^\pm$ is the Jordan decomposition (Theorem~\ref{signed-measures:theorem-jordan}) of the signed measure $\mu^+ - \mu^-$; the partition formula for
$\|\mu\|$ above is Proposition~\ref{signed-measures:proposition-variation-signed}(2).

For a noncompact $X$ the natural Banach space is $C_0(X)$, the space of continuous $f \colon X \to \RR$ \emph{vanishing at infinity}: for every $\varepsilon > 0$ the set $\{|f| \geq \varepsilon\}$ is compact.
With the norm $\|f\|_\infty$ it is a closed subspace of the bounded continuous functions, hence complete (a uniform limit of functions vanishing at infinity vanishes at infinity), and it contains
$C_c(X)$ as a dense subspace: for $f \in C_0(X)$ and $\varepsilon > 0$, Lemma~\ref{general-topology:lemma-lch-shrinking} gives $\psi \prec X$ with $\psi = 1$ on the compact set $\{|f| \geq \varepsilon\}$, and
then $\psi f \in C_c(X)$ and $\|f - \psi f\|_\infty \leq \varepsilon$. For $X = \RR^n$, $e^{-|x|^2}$ and $1/(1 + |x|^2)$ lie in $C_0$, while the constant $1$ does not. As in Corollary~\ref{radon-measures:corollary-signed-riesz}, a
\emph{finite signed Radon measure} is a difference $\mu = \mu^+ - \mu^-$ of finite Radon measures concentrated on disjoint Borel sets, with total variation norm $\|\mu\| = \mu^+(X) + \mu^-(X)$ and
$\int f\,d\mu = \int f\,d\mu^+ - \int f\,d\mu^-$.

\begin{corollary}[The dual of $C_0(X)$]\label{radon-measures:corollary-riesz-c0}
For every bounded linear functional $L$ on $C_0(X)$ there is a unique finite signed Radon measure $\mu$ with
\[
L(f) = \int f\,d\mu \quad (f \in C_0(X)), \qquad \|L\| = \|\mu\|.
\]
Conversely every finite signed Radon measure defines in this way a bounded functional of norm $\|\mu\|$. So $C_0(X)^*$ is isometrically isomorphic to the space of finite signed Radon measures with the
total variation norm. More generally, with $H$ as above, every bounded linear functional on $C_0(X, H)$ (with the norm $\sup_x|f(x)|$) is $f \mapsto \int\langle f, \nu\rangle\,d\mu$ with $\mu$ a finite Radon
measure of mass $\mu(X) = \|L\|$ and $|\nu| = 1$ $\mu$-almost everywhere.
\end{corollary}

\begin{proof}
We prove the vector-valued statement; the scalar one is the case $H = \RR$, with $\mu^\pm$ as in the proof of Corollary~\ref{radon-measures:corollary-signed-riesz}. The restriction of $L$ to $C_c(X, H)$
satisfies the hypothesis of Theorem~\ref{radon-measures:theorem-riesz-vector} with $M_K \leq \|L\|$, which gives $\mu$ and $\nu$. By \eqref{radon-measures:equation-total-variation} with $V = X$,
$\mu(X) = \sup\{L(f) : f \in C_c(X, H), |f| \leq 1\} \leq \|L\|$; in particular $\mu$ is finite. The functional $f \mapsto \int\langle f, \nu\rangle\,d\mu$ is defined on $C_0(X, H)$, with
$|\int\langle f, \nu\rangle\,d\mu| \leq \|f\|_\infty\mu(X)$, and it agrees with $L$ on $C_c(X, H)$. The density argument above, applied componentwise (or with $\psi f$ for $f \in C_0(X, H)$), shows that
$C_c(X, H)$ is dense in $C_0(X, H)$; two bounded functionals agreeing on a dense subspace are equal, so $L(f) = \int\langle f, \nu\rangle\,d\mu$ on $C_0(X, H)$. This also gives $\|L\| \leq \mu(X)$, so
$\|L\| = \mu(X)$. Uniqueness follows from the uniqueness in Theorem~\ref{radon-measures:theorem-riesz-vector}, applied to the restriction to $C_c(X, H)$. For the converse, let $\mu = \mu^+ - \mu^-$ with $\mu^\pm$ concentrated on disjoint Borel sets
$X^\pm$. The functional $L(f) = \int f\,d\mu$ satisfies $|L(f)| \leq \|f\|_\infty\|\mu\|$, and it is represented as in Theorem~\ref{radon-measures:theorem-riesz-vector} by the pair
$(\mu^+ + \mu^-, \mathbf{1}_{X^+} - \mathbf{1}_{X^-})$ (a finite Borel measure, hence Radon by Proposition~\ref{radon-measures:proposition-radon-automatic}). By uniqueness and the first part,
$\|L\| = (\mu^+ + \mu^-)(X) = \|\mu\|$.
\end{proof}

\begin{remark}\label{radon-measures:remark-dual-bounded-continuous}
For a noncompact $X$ the corollary does not extend from $C_0(X)$ to the space $C_b(X)$ of all bounded continuous functions. For $X = \NN$ with the discrete topology, $C_0(\NN) = c_0$ is the space of
sequences tending to $0$, its dual is the space of finite signed measures on $\NN$, that is, of summable sequences $\ell^1$, and $C_b(\NN) = \ell^\infty$. A bounded linear functional $L$ on $\ell^\infty$ that
vanishes on $c_0$ and has $L(\mathbf{1}) = 1$ cannot be integration against a countably additive measure: such a measure would give mass $L(\mathbf{1}_{\{n\}}) = 0$ to each point, hence mass $0$ to $\NN$.
Such functionals exist, by the Hahn--Banach theorem (Chapter~\ref{functional-analysis}); they correspond to finitely additive measures, or to measures on a large compactification of $X$. This is why the
compactness of $X$ matters in Corollary~\ref{radon-measures:corollary-signed-riesz}, and why weak convergence of measures, in the next section, is defined by testing against $C_c(X)$; testing
against $C_b(X)$ is a stronger condition, which also requires that no mass escapes to infinity.
\end{remark}

\begin{example}\label{radon-measures:example-riesz-vector}
\begin{enumerate}
\item For $L(\varphi) = \varphi(0) - \varphi(1)$ on $C_c(\RR)$, $\mu = \delta_0 + \delta_1$, $\nu(0) = 1$ and $\nu(1) = -1$. Indeed, for open $V$, the supremum in
\eqref{radon-measures:equation-total-variation} is the number of the points $0, 1$ lying in $V$, attained by functions equal to $1$ near $0$ and $-1$ near $1$.
\item For $h \in L^1_{\mathrm{loc}}(\RR^n)$ and $L(\varphi) = \int\varphi h\,d\lambda$, $\mu = |h|\lambda$ and $\nu = h/|h|$ where $h \neq 0$ (and $\nu = 1$ elsewhere, a $\mu$-null set). The total variation of
$L$ on $V$ is $\int_V|h|\,d\lambda$, the $L^1$ norm.
\item For the flux $L(F) = \int_0^{2\pi}\langle F(\cos t, \sin t), (\cos t, \sin t)\rangle\,dt = \int\langle F, n\rangle\,d\sigma$, with $\sigma$ the arc length measure of
Example~\ref{radon-measures:example-positive-functionals}(3), $\mu = \sigma$ and $\nu = n$ on $S^1$: the theorem recovers the curve and its unit normal from the functional.
For a $C^1$ field $F$, the divergence theorem in the plane gives $L(F) = \int_{|x| < 1}\operatorname{div}F\,d\lambda^2$.
\end{enumerate}
\end{example}

\begin{remark}\label{radon-measures:remark-perimeter}
\emph{Sets of finite perimeter.} Example~\ref{radon-measures:example-riesz-vector}(3) suggests a definition of boundary measure for arbitrary sets. For a Lebesgue measurable $E \subseteq \RR^n$ consider
$L_E(F) = \int_E\operatorname{div}F\,d\lambda^n$ for $C^1$ vector fields $F$ with compact support. If $\sup\{L_E(F) : |F| \leq 1, \operatorname{supp}F \subseteq K\} < \infty$ for every compact $K$, then $L_E$ extends
by uniform approximation to $C_c(\RR^n, \RR^n)$, and Theorem~\ref{radon-measures:theorem-riesz-vector} gives a Radon measure $\mu_E$, the \emph{perimeter measure}, and a unit vector field $\nu_E$ with
$\int_E\operatorname{div}F = \int\langle F, \nu_E\rangle\,d\mu_E$: a divergence theorem for $E$, valid with no regularity of $\partial E$ assumed. For a bounded open set with smooth boundary, $\mu_E$ is the
$(n-1)$-dimensional Hausdorff measure $\mathcal{H}^{n-1}$ (Chapter~\ref{geometric-measure-theory}) on $\partial E$ and $\nu_E$ is the outer normal. Such sets $E$ are the sets of \emph{finite perimeter} of Caccioppoli and De Giorgi, and the same construction with $\int u\operatorname{div}F$
defines functions of \emph{bounded variation}; De Giorgi's structure theorem shows that $\mu_E$ is $\mathcal{H}^{n-1}$ restricted to a countably $(n-1)$-rectifiable ``reduced boundary''
\cite[Chapter 5]{evans-gariepy}. This is the setting in which the isoperimetric problem and minimal surfaces are solved in geometric measure theory.
\end{remark}

\section{Haar measure}\label{radon-measures:section-haar}

Lebesgue measure on $\RR^n$ is characterised, up to a constant factor, by being invariant under translations and finite on compact sets (Example~\ref{radon-measures:example-haar}(1) below). Many other groups
carry such a measure: the circle, with arc length; the multiplicative group $\RR_{>0}$, with $dx/x$; matrix groups; discrete groups, with counting measure. The Riesz representation theorem gives a uniform
construction for all locally compact groups. Haar's theorem is the basis of harmonic analysis on groups (Fourier analysis on $\RR^n$, Fourier series on the circle, Peter--Weyl theory on compact
groups) and of averaging arguments such as the unitarity of representations of compact groups. Throughout, $G$ is a topological group (Definition~\ref{topological-manifolds:definition-topological-group})
which is locally compact, Hausdorff, and has a countable basis.

\begin{notation}\label{radon-measures:notation-translates}
For $E \subseteq G$ and $x \in G$ write $xE = \{xy : y \in E\}$, $Ex = \{yx : y \in E\}$ and $E^{-1} = \{y^{-1} : y \in E\}$. For a function $f$ on $G$ and $y \in G$, the \emph{left translate} is
$L_yf(x) = f(y^{-1}x)$ and the \emph{right translate} is $R_yf(x) = f(xy)$; the inverses are chosen so that $L_{yz} = L_yL_z$ and $R_{yz} = R_yR_z$, and $L_y\mathbf{1}_E = \mathbf{1}_{yE}$. We write
$C_c^+(G)$ for the set of $f \in C_c(G)$ with $f \geq 0$ and $f \neq 0$. The group operation is written multiplicatively, with identity $e$, even for abelian groups such as $\RR^n$.
\end{notation}

\begin{definition}\label{radon-measures:definition-haar}
A \emph{left Haar measure} on $G$ is a nonzero Borel measure $\mu$, finite on compact sets, that is \emph{left invariant}: $\mu(xE) = \mu(E)$ for all $x \in G$ and Borel $E \subseteq G$. A \emph{right Haar
measure} is defined in the same way with $\mu(Ex) = \mu(E)$.
\end{definition}

By Proposition~\ref{radon-measures:proposition-radon-automatic} a Haar measure is a Radon measure. Left invariance for sets is equivalent to $\int L_yf\,d\mu = \int f\,d\mu$, that is,
$\int f(y^{-1}x)\,d\mu(x) = \int f(x)\,d\mu(x)$, for all $y$ and all $f \in C_c(G)$ (or all Borel $f \geq 0$). Indeed, the formula for $f = \mathbf{1}_E$ is the invariance of sets, it extends to simple and
then to nonnegative functions by monotone convergence; conversely, if it holds on $C_c(G)$, the measures $E \mapsto \mu(yE)$ and $\mu$ represent the same functional, and agree by the uniqueness in the
Riesz--Markov theorem (Theorem~\ref{radon-measures:theorem-riesz-markov}).

\begin{example}\label{radon-measures:example-haar}
\begin{enumerate}
\item On $\RR^n$, Lebesgue measure is a Haar measure (Theorem~\ref{measure-theory:theorem-lebesgue-measure}). On a discrete group, such as $\ZZ^n$ or a finite group, counting measure is a left and right Haar
measure: compact sets are finite.
\item On $\RR_{>0}$ under multiplication, $\mu(E) = \int_E\frac{dx}{x}$ is a Haar measure: $\mu(aE) = \int_{aE}\frac{dx}{x} = \int_E\frac{a\,dt}{at} = \mu(E)$ by the substitution $x = at$. It is the image of
Lebesgue measure under the isomorphism $\exp \colon (\RR, +) \to (\RR_{>0}, \cdot)$, and $[1, 2]$ and $[10, 20]$ have the same measure $\log 2$. On $\RR^\times = \RR \setminus \{0\}$ the Haar
measure is $dx/|x|$, and on $\CC^\times$ it is $dx\,dy/(x^2 + y^2)$.
\item On the circle $S^1 = \{z \in \CC : |z| = 1\}$, the image of Lebesgue measure on $[0, 2\pi)$ under $\theta \mapsto e^{i\theta}$ is a Haar measure (Example~\ref{measure-theory:example-image-measure}(3)).
\item On $\GL_n(\RR)$, an open subset of $\RR^{n^2}$, the measure $\mu(E) = \int_E|\det X|^{-n}\,dX$, with $dX$ Lebesgue measure on $\RR^{n^2}$, is a left and right Haar measure. For
$A \in \GL_n(\RR)$, the linear map $X \mapsto AX$ of $\RR^{n^2}$ acts on each of the $n$ columns of $X$ by $A$, so its determinant is $(\det A)^n$, and by the change of variables formula
(Theorem~\ref{measure-theory:theorem-change-of-variables}, or Lemma~\ref{measure-theory:lemma-linear-change}),
\[
\mu(AE) = \int_{AE}|\det X|^{-n}\,dX = \int_E|\det AY|^{-n}|\det A|^n\,dY = \mu(E).
\]
The same computation with rows applies to $X \mapsto XA$.
\end{enumerate}
\end{example}

\subsection{Existence}\label{radon-measures:subsection-haar-existence}

The construction imitates the definition of Lebesgue measure by covers: to measure a function $f$, count how many small translates of a fixed bump $\varphi$ are needed to dominate $f$, compare with a
fixed function $f_0$, and let the support of $\varphi$ shrink to the identity. Two lemmas prepare this.

\begin{lemma}\label{radon-measures:lemma-uniform-continuity}
Let $f \in C_c(G)$ and $\varepsilon > 0$. There is a neighbourhood $V$ of $e$ such that $|f(yx) - f(x)| < \varepsilon$ and $|f(xy) - f(x)| < \varepsilon$ for all $x \in G$ and $y \in V$.
\end{lemma}

\begin{proof}
We prove the first inequality; the second follows by applying the first to the function $x \mapsto f(x^{-1})$, which is in $C_c(G)$ because inversion is a homeomorphism. Let $K = \operatorname{supp}f$.
For each $x \in K$ continuity of $y \mapsto f(yx)$ at $y = e$ gives an open neighbourhood $U_x$ of $e$ with $|f(yx) - f(x)| < \varepsilon/2$ for $y \in U_x$, and continuity of multiplication gives an
open neighbourhood $W_x$ of $e$ with $W_x = W_x^{-1}$ and $W_xW_x \subseteq U_x$. The open sets $W_xx$ cover $K$; choose $x_1, \ldots, x_m$ with $K \subseteq \bigcup_jW_{x_j}x_j$, and let
$V = \bigcap_jW_{x_j}$, a neighbourhood of $e$ with $V = V^{-1}$.

Let $y \in V$. If $x \in K$, say $x \in W_{x_j}x_j$, then $xx_j^{-1} \in W_{x_j} \subseteq U_{x_j}$ and $yxx_j^{-1} \in VW_{x_j} \subseteq U_{x_j}$, so $|f(x) - f(x_j)| < \varepsilon/2$ and
$|f(yx) - f(x_j)| < \varepsilon/2$, whence $|f(yx) - f(x)| < \varepsilon$. If $x \notin K$ but $yx \in K$, the previous case applied to the point $yx$ and to $y^{-1} \in V$ gives
$|f(x) - f(yx)| < \varepsilon$. If neither $x$ nor $yx$ lies in $K$, both values vanish.
\end{proof}

\begin{lemma}\label{radon-measures:lemma-haar-positive}
Let $\mu$ be a left Haar measure. Then $\mu(U) > 0$ for every nonempty open $U$, and $\int f\,d\mu > 0$ for every $f \in C_c^+(G)$.
\end{lemma}

\begin{proof}
If $\mu(U) = 0$, every compact $K$ is covered by finitely many translates $xU$, each of measure $0$, so $\mu(K) = 0$; by inner regularity $\mu = 0$, which is excluded. If $f \in C_c^+(G)$, then
$U = \{f > \frac12\max f\}$ is a nonempty open set and $\int f\,d\mu \geq \frac12\max f\cdot\mu(U) > 0$.
\end{proof}

\begin{theorem}[Haar]\label{radon-measures:theorem-haar}
Every locally compact Hausdorff group $G$ with a countable basis has a left Haar measure. It is unique up to a positive factor: if $\mu$ and $\nu$ are left Haar measures, then $\nu = c\mu$ for some
$c > 0$.
\end{theorem}

\begin{proof}
\emph{Step 1: covering numbers.} For $f, \varphi \in C_c^+(G)$ let
\[
(f : \varphi) = \inf\Bigl\{\sum_{j=1}^mc_j : c_j \geq 0, \ y_j \in G, \ f \leq \sum_{j=1}^mc_jL_{y_j}\varphi\Bigr\},
\]
the cost of covering $f$ by translates of $\varphi$. It is finite: $\varphi > a$ on a nonempty open set $U$ for some $a > 0$, so $L_y\varphi > a$ on $yU$; finitely many $y_jU$ cover the compact support of $f$,
and $f \leq \sum_j(\max f/a)L_{y_j}\varphi$. Directly from the definition, for $f, g, \varphi, \psi \in C_c^+(G)$, $c > 0$ and $y \in G$:
\begin{enumerate}
\item $(L_yf : \varphi) = (f : \varphi)$, $(cf : \varphi) = c(f : \varphi)$, and $(f : \varphi) \leq (g : \varphi)$ if $f \leq g$;
\item $(f + g : \varphi) \leq (f : \varphi) + (g : \varphi)$;
\item $(f : \varphi) \leq (f : \psi)(\psi : \varphi)$: if $f \leq \sum_jc_jL_{y_j}\psi$ and $\psi \leq \sum_kd_kL_{z_k}\varphi$, then $f \leq \sum_{j,k}c_jd_kL_{y_jz_k}\varphi$;
\item $(f : \varphi) \geq \max f/\max\varphi > 0$, by evaluating a covering inequality at a point where $f$ is maximal.
\end{enumerate}
Fix $f_0 \in C_c^+(G)$ and put $I_\varphi(f) = (f : \varphi)/(f_0 : \varphi)$. By (3), applied to $(f_0 : \varphi) \leq (f_0 : f)(f : \varphi)$ and to $(f : \varphi) \leq (f : f_0)(f_0 : \varphi)$,
\begin{equation}\label{radon-measures:equation-haar-bounds}
\frac{1}{(f_0 : f)} \leq I_\varphi(f) \leq (f : f_0).
\end{equation}
Each $I_\varphi$ is left invariant, positively homogeneous, monotone and subadditive, by (1) and (2).

\emph{Step 2: approximate additivity.} We show: for $f_1, f_2 \in C_c^+(G)$ and $\varepsilon > 0$ there is a neighbourhood $V$ of $e$ such that
$I_\varphi(f_1) + I_\varphi(f_2) \leq I_\varphi(f_1 + f_2) + \varepsilon$ whenever $\operatorname{supp}\varphi \subseteq V$. Choose $g \in C_c(G)$ with $0 \leq g \leq 1$ and $g = 1$ on
$\operatorname{supp}(f_1 + f_2)$ (Lemma~\ref{general-topology:lemma-lch-shrinking}), let $\delta > 0$, and put $h = f_1 + f_2 + \delta g$ and $h_i = f_i/h$, with $h_i = 0$ where $f_i = 0$. Since $h \geq \delta/2$ on
the open set $\{g > \frac12\}$, which contains $\operatorname{supp}f_i$, the functions $h_i$ are continuous with compact support, and $h_1 + h_2 \leq 1$. By
Lemma~\ref{radon-measures:lemma-uniform-continuity} there is a neighbourhood $V$ of $e$ with $|h_i(yu) - h_i(y)| < \delta$ for $u \in V$, $y \in G$, $i = 1, 2$. Let $\operatorname{supp}\varphi \subseteq V$ and
$h \leq \sum_jc_jL_{y_j}\varphi$. Where $L_{y_j}\varphi(x) = \varphi(y_j^{-1}x) \neq 0$ we have $x = y_ju$ with $u \in V$, so $h_i(x) \leq h_i(y_j) + \delta$; hence
\[
f_i = h_ih \leq \sum_jc_jh_i(x)L_{y_j}\varphi(x) \leq \sum_jc_j\bigl(h_i(y_j) + \delta\bigr)L_{y_j}\varphi(x),
\]
so $(f_1 : \varphi) + (f_2 : \varphi) \leq \sum_jc_j(h_1(y_j) + h_2(y_j) + 2\delta) \leq (1 + 2\delta)\sum_jc_j$. Taking the infimum over the covers of $h$ and using (1), (2),
$(f_1 : \varphi) + (f_2 : \varphi) \leq (1 + 2\delta)\bigl((f_1 + f_2 : \varphi) + \delta(g : \varphi)\bigr)$. Dividing by $(f_0 : \varphi)$ and using \eqref{radon-measures:equation-haar-bounds},
\[
I_\varphi(f_1) + I_\varphi(f_2) \leq I_\varphi(f_1 + f_2) + 2\delta(f_1 + f_2 : f_0) + \delta(1 + 2\delta)(g : f_0),
\]
and the error is less than $\varepsilon$ for $\delta$ small enough; $V$ depends on $\delta$, but not on $\varphi$.

\emph{Step 3: a limit by compactness.} The product $P = \prod_{f \in C_c^+(G)}[1/(f_0 : f), (f : f_0)]$ of compact intervals is compact (Theorem~\ref{general-topology:theorem-tychonoff}), and by
\eqref{radon-measures:equation-haar-bounds} each $I_\varphi$ is a point of $P$. For a neighbourhood $V$ of $e$ let $\Phi_V$ be the closure in $P$ of $\{I_\varphi : \varphi \in C_c^+(G),\
\operatorname{supp}\varphi \subseteq V\}$, which is nonempty by Lemma~\ref{general-topology:lemma-lch-shrinking}. Since $\Phi_{V_1} \cap \cdots \cap \Phi_{V_m} \supseteq \Phi_{V_1 \cap \cdots \cap V_m} \neq \emptyset$,
the closed sets $\Phi_V$ have the finite intersection property, and by compactness there is a point $I \in \bigcap_V\Phi_V$. Concretely: for every neighbourhood $V$ of $e$, every finite set
$f_1, \ldots, f_k \in C_c^+(G)$ and every $\eta > 0$ there is $\varphi$ with $\operatorname{supp}\varphi \subseteq V$ and $|I(f_i) - I_\varphi(f_i)| < \eta$ for $i \leq k$.

\emph{Step 4: properties of $I$.} Any identity or inequality involving finitely many values of $I_\varphi$ and holding for all $\varphi$ passes to $I$ by the last remark. So $I$ is left invariant,
positively homogeneous and subadditive on $C_c^+(G)$, with $I(f) \geq 1/(f_0 : f) > 0$. Given $f_1, f_2$ and $\varepsilon$, take $V$ as in Step 2 and $\varphi$ with support in $V$ approximating $I$ at
$f_1$, $f_2$, $f_1 + f_2$ within $\varepsilon$: then $I(f_1) + I(f_2) \leq I(f_1 + f_2) + 4\varepsilon$. Hence $I$ is additive on $C_c^+(G)$.

\emph{Step 5: the measure.} Put $I(0) = 0$. For $f \in C_c(G)$ write $f = f^+ - f^-$ and set $I(f) = I(f^+) - I(f^-)$; if $f = g - h$ with $g, h \geq 0$ in $C_c(G)$, then $g + f^- = h + f^+$, and additivity gives
$I(f) = I(g) - I(h)$. This makes $I$ a positive linear functional on $C_c(G)$, and the Riesz--Markov theorem (Theorem~\ref{radon-measures:theorem-riesz-markov}) gives a Radon measure $\mu$ with
$\int f\,d\mu = I(f)$. It is nonzero, as $I(f_0) = 1$, and left invariant, since $\int L_yf\,d\mu = I(L_yf) = I(f)$ for $f \in C_c(G)$.

\emph{Step 6: uniqueness.} Let $\mu$, $\nu$ be left Haar measures. $G$ is $\sigma$-compact (it has an exhaustion by compact sets), so both are $\sigma$-finite, and
$\mathcal{B}(G \times G) = \mathcal{B}(G) \otimes \mathcal{B}(G)$ (Proposition~\ref{measure-theory:proposition-borel-product}); so Tonelli's theorem applies to nonnegative continuous functions on $G \times G$. Fix
$f \in C_c^+(G)$ and put $F(x) = \int f(tx)\,d\mu(t)$. Then $F(x) > 0$ by Lemma~\ref{radon-measures:lemma-haar-positive}, applied to $t \mapsto f(tx)$, and $F$ is continuous: given $x$ and $\varepsilon$,
Lemma~\ref{radon-measures:lemma-uniform-continuity} gives a neighbourhood $V$ of $e$ with compact closure and $|f(txu) - f(tx)| < \varepsilon$ for all $t$ and $u \in V$; both terms vanish unless
$t \in C = \operatorname{supp}(f)\overline{V}^{-1}x^{-1}$, a compact set, so $|F(xu) - F(x)| \leq \varepsilon\mu(C)$. For $g \in C_c^+(G)$, the function $(x, y) \mapsto g(x)f(yx)/F(x)$ is continuous and
nonnegative, and we integrate it in two orders. First,
\[
\int\!\!\int\frac{g(x)f(yx)}{F(x)}\,d\mu(y)\,d\nu(x) = \int g(x)\frac{F(x)}{F(x)}\,d\nu(x) = \int g\,d\nu.
\]
Second, using left invariance of $\nu$ ($x \mapsto y^{-1}x$) and then of $\mu$ ($y \mapsto xy$),
\[\begin{aligned}
\int\!\!\int\frac{g(x)f(yx)}{F(x)}\,d\nu(x)\,d\mu(y) &= \int\!\!\int\frac{g(y^{-1}x)f(x)}{F(y^{-1}x)}\,d\nu(x)\,d\mu(y) \\
&= \int f(x)\int\frac{g(y^{-1}x)}{F(y^{-1}x)}\,d\mu(y)\,d\nu(x) \\
&= \int f(x)\int\frac{g(y^{-1})}{F(y^{-1})}\,d\mu(y)\,d\nu(x) = c_g\int f\,d\nu,
\end{aligned}\]
where $c_g = \int g(y^{-1})/F(y^{-1})\,d\mu(y)$ depends on $\mu$, $f$, $g$ but not on $\nu$. Hence $\int g\,d\nu = c_g\int f\,d\nu$, and with $\nu = \mu$, $\int g\,d\mu = c_g\int f\,d\mu$. Therefore
$\int g\,d\nu = c\int g\,d\mu$ for all $g \in C_c^+(G)$, with $c = \int f\,d\nu/\int f\,d\mu > 0$ (Lemma~\ref{radon-measures:lemma-haar-positive}); by linearity this holds on $C_c(G)$, and $\nu = c\mu$ by
the uniqueness in the Riesz--Markov theorem.
\end{proof}

\begin{remark}\label{radon-measures:remark-haar}
The theorem holds for all locally compact Hausdorff groups, without the countable basis, with the same proof once the Riesz representation theorem and Fubini's theorem are available for Radon
measures in that generality \cite[Section 11.1]{folland-real}. The first construction, for compact metric groups, is due to Haar (1933); von Neumann proved uniqueness, and Weil gave the general
case. For Lie groups an invariant measure can be written down as the integral of a left-invariant volume form (Chapter~\ref{lie-groups}).
\end{remark}

\subsection{The modular function}\label{radon-measures:subsection-modular}

Left Haar measure need not be right invariant, but right translation can only change it by a factor.

\begin{proposition}\label{radon-measures:proposition-modular}
Let $\mu$ be a left Haar measure on $G$. There is a unique map $\Delta \colon G \to \RR_{>0}$, the \emph{modular function}, with $\mu(Ex) = \Delta(x)\mu(E)$ for all $x \in G$ and Borel $E$. It does not depend on
the choice of $\mu$, it is a continuous homomorphism, and $\int R_xf\,d\mu = \Delta(x^{-1})\int f\,d\mu$ for $f \in C_c(G)$.
\end{proposition}

\begin{proof}
For fixed $x$, $E \mapsto \mu(Ex)$ is a Borel measure, finite on compact sets (as $Kx$ is compact), nonzero, and left invariant, since $(yE)x = y(Ex)$. By uniqueness it is $\Delta(x)\mu$ for a number
$\Delta(x) > 0$; replacing $\mu$ by $c\mu$ does not change $\Delta(x)$. From $\mu(Exy) = \Delta(y)\mu(Ex) = \Delta(y)\Delta(x)\mu(E)$ we get $\Delta(xy) = \Delta(x)\Delta(y)$. For $f = \mathbf{1}_E$,
$R_x\mathbf{1}_E = \mathbf{1}_{Ex^{-1}}$, so $\int R_x\mathbf{1}_E\,d\mu = \Delta(x^{-1})\mu(E)$; this extends to $C_c(G)$ as before. Finally, for $f \in C_c^+(G)$, $\Delta(x^{-1}) = \int f(tx)\,d\mu(t)/\int f\,d\mu$
is continuous in $x$, by the argument for $F$ in Step 6 of the proof of Theorem~\ref{radon-measures:theorem-haar}, and so is its inverse $\Delta$.
\end{proof}

$G$ is \emph{unimodular} if $\Delta = 1$, that is, if a left Haar measure is also right invariant.

\begin{corollary}\label{radon-measures:corollary-unimodular}
Abelian groups, discrete groups and compact groups are unimodular. On a unimodular group, a Haar measure is invariant under inversion: $\mu(E^{-1}) = \mu(E)$. On a compact group there is a unique Haar
probability measure, and it is left and right invariant and invariant under inversion.
\end{corollary}

\begin{proof}
For abelian groups $Ex = xE$. For discrete groups counting measure is right invariant. If $G$ is compact, $\Delta(G)$ is a compact subgroup of $(\RR_{>0}, \cdot)$, hence $\{1\}$: a subgroup containing some
$a \neq 1$ contains all $a^n$, $n \in \ZZ$, which is unbounded or accumulates at $0$. If $G$ is unimodular, $\check\mu(E) = \mu(E^{-1})$ is a Borel measure, finite on compacts, and left invariant,
since $(xE)^{-1} = E^{-1}x^{-1}$ and $\mu$ is right invariant; so $\check\mu = c\mu$, and applying this twice gives $c^2 = 1$, $c = 1$. On a compact group $0 < \mu(G) < \infty$, and $\mu/\mu(G)$ is the
unique left Haar measure of total mass $1$.
\end{proof}

\begin{example}\label{radon-measures:example-ax-plus-b}
\emph{A group that is not unimodular.} Let $G = \{(a, b) : a > 0, b \in \RR\}$ with $(a, b)(a^{\prime}, b^{\prime}) = (aa^{\prime}, ab^{\prime} + b)$, the group of affine maps $t \mapsto at + b$ of the line, an open subset
of $\RR^2$. Left multiplication by $(a, b)$ is the map $(a^{\prime}, b^{\prime}) \mapsto (aa^{\prime}, ab^{\prime} + b)$, with Jacobian determinant $a^2$; by the change of variables formula
$d\mu = a^{-2}\,da\,db$ is a left Haar measure. Right multiplication by $(a^{\prime}, b^{\prime})$ is $(a, b) \mapsto (aa^{\prime}, ab^{\prime} + b)$, with Jacobian determinant $a^{\prime}$, so $\mu(E(a^{\prime}, b^{\prime})) = \mu(E)/a^{\prime}$:
$\Delta(a, b) = 1/a$, and $a^{-1}\,da\,db$ is a right Haar measure, different from the left one.
\end{example}

\subsection{Steinhaus' theorem}\label{radon-measures:subsection-steinhaus}

\begin{theorem}[Steinhaus]\label{radon-measures:theorem-steinhaus}
Let $\mu$ be a left Haar measure on $G$ and $A \subseteq G$ a Borel set with $\mu(A) > 0$. Then $AA^{-1} = \{ab^{-1} : a, b \in A\}$ contains a neighbourhood of $e$.
\end{theorem}

\begin{proof}
By inner regularity choose a compact $K \subseteq A$ with $\mu(K) > 0$, and by outer regularity an open $U \supseteq K$ with $\mu(U) < 2\mu(K)$. The continuous map $(v, k) \mapsto vk$ sends $\{e\} \times K$ into $U$,
so each $(e, k)$ has a product neighbourhood $V_k \times W_k$ mapped into $U$; finitely many $W_{k_i}$ cover $K$, and $V = \bigcap_iV_{k_i}$ is a neighbourhood of $e$ with $VK \subseteq U$. For $v \in V$ the
sets $vK$ and $K$ lie in $U$ and have measure $\mu(K)$; if they were disjoint, $\mu(U) \geq 2\mu(K)$. So $vk = k^{\prime}$ for some $k, k^{\prime} \in K$, and $v = k^{\prime}k^{-1} \in KK^{-1} \subseteq AA^{-1}$.
\end{proof}

\begin{corollary}\label{radon-measures:corollary-measurable-homomorphism}
A Borel measurable homomorphism $\phi \colon (\RR, +) \to (\RR, +)$ is of the form $\phi(x) = cx$. In particular every Lebesgue measurable solution of Cauchy's equation $\phi(x + y) = \phi(x) + \phi(y)$
is linear.
\end{corollary}

\begin{proof}
For a Lebesgue measurable $\phi$, the sets $A_n = \phi^{-1}((-n, n))$ are Lebesgue measurable and cover $\RR$, so some $A_n$ has positive measure, and so does a Borel subset of it (the measurable sets are the
Borel sets up to null sets, Proposition~\ref{measure-theory:proposition-lebesgue-completion}). By Steinhaus' theorem $A_n - A_n \supseteq (-\delta, \delta)$ for some $\delta > 0$, and
$\phi(a - b) = \phi(a) - \phi(b) \in (-2n, 2n)$; so $|\phi| < 2n$ on $(-\delta, \delta)$. Additivity gives $\phi(x/m) = \phi(x)/m$, so $|\phi| < 2n/m$ on $(-\delta/m, \delta/m)$: $\phi$ is continuous at
$0$, hence everywhere, as $\phi(x + h) - \phi(x) = \phi(h)$. Finally $\phi(q) = q\phi(1)$ for rational $q$ by additivity, and by continuity $\phi(x) = \phi(1)x$.
\end{proof}

Without measurability the statement fails: viewing $\RR$ as a vector space over $\QQ$ and choosing a basis (a \emph{Hamel basis}, which exists by the axiom of choice), any $\QQ$-linear map that is not
$\RR$-linear, such as a projection onto the line spanned by one basis vector, is a non-measurable solution of Cauchy's equation.

\section{Weak convergence of measures}\label{radon-measures:section-weak-convergence}

Sequences of measures appear whenever something discrete approximates something continuous, or something smooth approximates something singular. Riemann sums are integrals against finitely many
point masses that approximate Lebesgue measure; in probability the distributions of normalised sums of random variables approximate the Gaussian distribution; smooth bumps of total mass $1$
concentrating at a point approximate the Dirac measure; and in geometric measure theory one studies limits of the area measures of a sequence of surfaces, for instance of a minimising sequence for
the area. Convergence $\mu_k(A) \to \mu(A)$ for all Borel sets $A$ is far too strong for these purposes (Example~\ref{radon-measures:example-weak-convergence}(1)). The useful notion tests the measures
against continuous functions. As in Section~\ref{radon-measures:section-radon}, $X$ is a locally compact Hausdorff space with a countable basis.

\begin{definition}\label{radon-measures:definition-weak-convergence}
A sequence of Radon measures $\mu_k$ on $X$ \emph{converges weakly} to a Radon measure $\mu$, written $\mu_k \rightharpoonup \mu$, if
\[
\int_Xf\,d\mu_k \to \int_Xf\,d\mu \qquad \text{for every } f \in C_c(X).
\]
\end{definition}

The limit is unique, by the uniqueness in Theorem~\ref{radon-measures:theorem-riesz-markov}. In probability theory this is called \emph{vague} convergence, and \emph{weak} convergence of probability
measures means testing against all bounded continuous functions; we call it \emph{narrow} convergence, and the two agree when no mass is lost
(Lemma~\ref{radon-measures:lemma-weak-to-narrow}). In functional analysis it is the weak-$*$
convergence of functionals on $C_c(X)$.

\begin{example}\label{radon-measures:example-weak-convergence}
\begin{enumerate}
\item \emph{Moving point masses.} If $x_k \to x$ in $X$, then $\delta_{x_k} \rightharpoonup \delta_x$, since $f(x_k) \to f(x)$ for continuous $f$. On $\RR$, $\delta_{1/k} \rightharpoonup \delta_0$, but
$\delta_{1/k}(\{0\}) = 0$ does not converge to $\delta_0(\{0\}) = 1$, and $\delta_{1/k}((0, 1)) = 1$ does not converge to $\delta_0((0, 1)) = 0$. So weak convergence does not imply convergence on
compact sets, nor on open sets; Proposition~\ref{radon-measures:proposition-weak-convergence-sets} gives the inequalities that do hold.
\item \emph{Riemann sums.} On $\RR$ let $\mu_k = \frac1k\sum_{j=1}^k\delta_{j/k}$. Then $\mu_k \rightharpoonup \mathbf{1}_{[0, 1]}\lambda$. Indeed $f \in C_c(\RR)$ is uniformly continuous, being
continuous with compact support; with $\omega(\delta) = \sup\{|f(s) - f(t)| : |s - t| \leq \delta\}$, which tends to $0$ with $\delta$,
\[
\Bigl|\int f\,d\mu_k - \int_0^1f\,d\lambda\Bigr| = \Bigl|\sum_{j=1}^k\int_{(j-1)/k}^{j/k}\bigl(f(j/k) - f(t)\bigr)\,dt\Bigr| \leq \omega(1/k) \to 0.
\]
Each $\mu_k$ lives on a finite set, of Lebesgue measure $0$, while the limit is absolutely continuous with respect to $\lambda$.
\item \emph{Approximate identities.} Let $\phi \geq 0$ be continuous on $\RR$ with support in $[-1, 1]$ and $\int\phi\,d\lambda = 1$, and $\phi_k(x) = k\phi(kx)$. Then $\int\phi_k\,d\lambda = 1$
(Theorem~\ref{measure-theory:theorem-change-of-variables}), $\phi_k$ vanishes outside $[-1/k, 1/k]$, and for $f \in C_c(\RR)$
\[
\Bigl|\int f\phi_k\,d\lambda - f(0)\Bigr| = \Bigl|\int\bigl(f(x) - f(0)\bigr)\phi_k(x)\,dx\Bigr| \leq \omega(1/k),
\]
so $\phi_k\lambda \rightharpoonup \delta_0$: absolutely continuous measures converge to a singular one.
\item \emph{Loss of mass.} On $\RR$, $\delta_k \rightharpoonup 0$, since every $f \in C_c(\RR)$ vanishes at $k$ for $k$ large; and $\frac1k\mathbf{1}_{[0, k]}\lambda \rightharpoonup 0$, since
$|\int_0^kf\,d\lambda|/k \leq \|f\|_1/k$. All these measures have total mass $1$, the limit has mass $0$: mass can escape to infinity, by moving away or by spreading out.
\item \emph{Oscillation.} Let $g \geq 0$ be continuous on $\RR$ and periodic of period $2\pi$, with mean value $\bar g = \frac1{2\pi}\int_0^{2\pi}g\,d\lambda$. Then $g(kx)\lambda \rightharpoonup \bar g\lambda$; for
instance $(1 + \sin kx)\lambda \rightharpoonup \lambda$, although the densities $1 + \sin kx$ do not converge pointwise (at $x = \pi/2$ they run through $2, 1, 0, 1, 2, \ldots$). To see this, let
$f \in C_c(\RR)$ with support in $[a, b]$ and $\varepsilon > 0$, and choose $a = a_0 < \cdots < a_N = b$ with $|f(s) - f(a_j)| \leq \varepsilon$ for $s \in [a_{j-1}, a_j]$. An interval
$[k\alpha, k\beta]$ consists of $\lfloor k(\beta - \alpha)/2\pi\rfloor$ full periods and a remainder shorter than $2\pi$, so
$\int_\alpha^\beta g(kx)\,dx = \frac1k\int_{k\alpha}^{k\beta}g(t)\,dt = (\beta - \alpha)\bar g + r$ with $|r| \leq 4\pi\|g\|_\infty/k$. Replacing $f$ by $f(a_j)$ on $[a_{j-1}, a_j]$ costs at most
$\varepsilon(b - a)\|g\|_\infty$ in $\int fg(kx)\,dx$ and in $\bar g\int f\,dx$, so
\[
\Bigl|\int f(x)g(kx)\,dx - \bar g\int f(x)\,dx\Bigr| \leq 2\varepsilon(b - a)\|g\|_\infty + \frac{4\pi N\|f\|_\infty\|g\|_\infty}{k},
\]
whose $\limsup$ as $k \to \infty$ is at most $2\varepsilon(b - a)\|g\|_\infty$. Fast oscillations average out; this is the measure-theoretic content of the Riemann--Lebesgue lemma.
\end{enumerate}
\end{example}

\begin{proposition}\label{radon-measures:proposition-weak-convergence-sets}
Let $\mu_k \rightharpoonup \mu$ be Radon measures on $X$.
\begin{enumerate}
\item $\mu(U) \leq \liminf_k\mu_k(U)$ for every open $U$.
\item $\mu(K) \geq \limsup_k\mu_k(K)$ for every compact $K$.
\item $\mu_k(A) \to \mu(A)$ for every Borel set $A$ with compact closure and $\mu(\partial A) = 0$.
\end{enumerate}
\end{proposition}

\begin{proof}
(1) By inner regularity, $\mu(U)$ is the supremum of $\mu(K)$ over compact $K \subseteq U$; for such $K$, Lemma~\ref{general-topology:lemma-lch-shrinking} gives $f \prec U$ with $f = 1$ on $K$, and
$\mu(K) \leq \int f\,d\mu \leq \mu(U)$. So $\mu(U) = \sup\{\int f\,d\mu : f \prec U\}$. For $f \prec U$, $\int f\,d\mu = \lim_k\int f\,d\mu_k \leq \liminf_k\mu_k(U)$.

(2) Let $V \supseteq K$ be open. By Lemma~\ref{general-topology:lemma-lch-shrinking} there is $f \prec V$ with $f = 1$ on $K$, and $\limsup_k\mu_k(K) \leq \lim_k\int f\,d\mu_k = \int f\,d\mu \leq \mu(V)$.
By outer regularity, the infimum of $\mu(V)$ over such $V$ is $\mu(K)$.

(3) The interior $\operatorname{int}A$ and the closure $\overline{A}$ differ by $\partial A$, so $\mu(\operatorname{int}A) = \mu(A) = \mu(\overline{A})$. By (1) and (2),
$\mu(A) = \mu(\operatorname{int}A) \leq \liminf_k\mu_k(\operatorname{int}A) \leq \liminf_k\mu_k(A) \leq \limsup_k\mu_k(A) \leq \limsup_k\mu_k(\overline{A}) \leq \mu(\overline{A}) = \mu(A)$.
\end{proof}

In Example~\ref{radon-measures:example-weak-convergence}(1) the inequalities are strict: $\delta_0((0, 1)) = 0 < 1$ and $\delta_0(\{0\}) = 1 > 0$; the sets $(0, 1)$ and $\{0\}$ have the point $0$, of
$\delta_0$-measure $1$, in their boundary. In (2) of the same example, (3) says that the proportion of the points $j/k$ lying in an interval $[\alpha, \beta] \subseteq [0, 1]$ tends to $\beta - \alpha$.

The main theorem about weak convergence is a compactness theorem: bounded mass on compact sets suffices for a convergent subsequence. Its proof needs a countable supply of test functions.

\begin{lemma}\label{radon-measures:lemma-cc-separable}
Let $K_1 \subseteq K_2 \subseteq \cdots$ be an exhaustion of $X$ by compact sets with $K_m \subseteq \operatorname{int}K_{m+1}$. There is a countable set $\mathcal{D} \subseteq C_c(X)$ such that for every
$m$, every $f \in C_c(X)$ with $\operatorname{supp}f \subseteq K_m$ and every $\varepsilon > 0$ there is $g \in \mathcal{D}$ with $\operatorname{supp}g \subseteq K_{m+1}$ and $\|f - g\|_\infty < \varepsilon$.
\end{lemma}

\begin{proof}
First, $C(L, \RR)$ has a countable dense subset for every compact metric space $(L, d)$. Indeed $L$ has a countable dense subset $\{x_1, x_2, \ldots\}$ (for each $j$, finitely many balls of radius $1/j$
cover $L$; take their centres). The functions $u_i(y) = d(x_i, y)$ separate points: if $y \neq z$ and $d(x_i, y) < d(y, z)/2$, then $d(x_i, z) > d(y, z)/2 > d(x_i, y)$. By the Stone--Weierstrass theorem
(Theorem~\ref{real-analysis:theorem-stone-weierstrass}) the polynomials in finitely many of the $u_i$ are dense in $C(L, \RR)$, and so are those with rational coefficients, which form a countable set, since
the $u_i$ are bounded and real coefficients can be approximated by rational ones.

Each $K_{m+1}$ is a compact metric space, $X$ being metrisable; let $\mathcal{D}_m \subseteq C(K_{m+1}, \RR)$ be countable and dense. By Lemma~\ref{general-topology:lemma-lch-shrinking} choose
$\chi_m \in C_c(X)$ with values in $[0, 1]$, $\chi_m = 1$ on $K_m$ and $\operatorname{supp}\chi_m \subseteq \operatorname{int}K_{m+1}$. For $u \in \mathcal{D}_m$ the function $\chi_mu$, extended by $0$
outside $K_{m+1}$, is continuous on $X$: it is continuous on the open set $\operatorname{int}K_{m+1}$ and vanishes on the open set $X \setminus \operatorname{supp}\chi_m$, and these two open sets cover $X$.
Let $\mathcal{D}$ be the set of all these functions, for all $m$. Given $f$ with support in $K_m$, choose $u \in \mathcal{D}_m$ with $|f - u| < \varepsilon$ on $K_{m+1}$; then $f = \chi_mf$, so
$|f - \chi_mu| = \chi_m|f - u| < \varepsilon$ on $K_{m+1}$, and both functions vanish outside $K_{m+1}$.
\end{proof}

\begin{theorem}[Weak compactness of Radon measures]\label{radon-measures:theorem-weak-compactness}
Let $(\mu_k)$ be a sequence of Radon measures on $X$ with $\sup_k\mu_k(K) < \infty$ for every compact $K \subseteq X$. Then a subsequence converges weakly to a Radon measure.
\end{theorem}

\begin{proof}
Let $(K_m)$ be an exhaustion as in Lemma~\ref{radon-measures:lemma-cc-separable}, $C_m = \sup_k\mu_k(K_m)$, and $\mathcal{D} = \{g_1, g_2, \ldots\}$ the countable set of the lemma. For each $i$ the
sequence $\int g_i\,d\mu_k$ is bounded, by $\|g_i\|_\infty$ times $\sup_k\mu_k(\operatorname{supp}g_i)$, so it has a convergent subsequence (Theorem~\ref{real-analysis:theorem-bolzano-weierstrass-real}).
Choose a subsequence along which $\int g_1\,d\mu_k$ converges, a further subsequence along which $\int g_2\,d\mu_k$ converges, and so on; the diagonal sequence, the $k$-th term of the $k$-th
subsequence, is a subsequence of each of them from its $i$-th term on. Renaming it $(\mu_k)$, $\int g\,d\mu_k$ converges for every $g \in \mathcal{D}$.

Let $f \in C_c(X)$, with support in some $K_m$, and $\varepsilon > 0$. Choose $g \in \mathcal{D}$ with support in $K_{m+1}$ and $\|f - g\|_\infty < \varepsilon$. Then
\[
\Bigl|\int f\,d\mu_k - \int f\,d\mu_l\Bigr| \leq 2\varepsilon C_{m+1} + \Bigl|\int g\,d\mu_k - \int g\,d\mu_l\Bigr|,
\]
so $\int f\,d\mu_k$ is a Cauchy sequence. Its limit $\Lambda f$ is linear and positive in $f$, so by Theorem~\ref{radon-measures:theorem-riesz-markov} it is $\int f\,d\mu$ for a Radon measure $\mu$, and
$\mu_k \rightharpoonup \mu$.
\end{proof}

\begin{corollary}\label{radon-measures:corollary-probability-compact}
Let $X$ be a compact metric space. Every sequence of Borel probability measures on $X$ has a subsequence converging weakly to a Borel probability measure.
\end{corollary}

\begin{proof}
A compact metric space is locally compact and Hausdorff, and it has a countable basis: the balls of radius $1/j$ centred at the points of finite sets of such balls covering $X$. Every Borel measure on
$X$ that is finite is a Radon measure (Proposition~\ref{radon-measures:proposition-radon-automatic}), and $\mu_k(X) = 1$ bounds the masses. Theorem~\ref{radon-measures:theorem-weak-compactness} gives a
weakly convergent subsequence, and its limit $\mu$ satisfies $\mu(X) = \int 1\,d\mu = \lim\int 1\,d\mu_k = 1$, as $1 \in C_c(X)$ for compact $X$.
\end{proof}

\begin{remark}\label{radon-measures:remark-weak-compactness}
\begin{enumerate}
\item The hypothesis cannot be dropped: $k\delta_0$ has no weakly convergent subsequence on $\RR$. On non-compact spaces the limit of probability measures may have smaller mass, by
Example~\ref{radon-measures:example-weak-convergence}(4); by Proposition~\ref{radon-measures:proposition-weak-convergence-sets}(1) applied to $U = X$, mass can only be lost, never gained.
\item Theorem~\ref{radon-measures:theorem-weak-compactness} is a case of the Banach--Alaoglu theorem: bounded sets of linear functionals are compact for the weak-$*$ topology
(Chapter~\ref{functional-analysis}).
\item In the calculus of variations and in geometric measure theory this is the basic source of limits. A minimising sequence of surfaces for the area, all contained in a fixed ball, has area
measures with bounded mass, and a subsequence of these converges weakly to a Radon measure; the hard part, the \emph{regularity theory}, is to show that the limit is again the area measure of a
surface, perhaps with singularities. For curves, Exercise~\ref{radon-measures:exercise-staircase} shows that length can jump in the limit.
\end{enumerate}
\end{remark}

\subsection{Narrow convergence and Prokhorov's theorem}\label{radon-measures:subsection-prokhorov}

In probability one tests probability measures against all bounded continuous functions, not only those with compact support; then the total mass $\int1\,d\mu_k = 1$ is preserved in the limit,
and escape of mass to infinity, as in Example~\ref{radon-measures:example-weak-convergence}(4), is excluded. We write $C_b(X)$ for the space of bounded continuous functions $X \to \RR$.

\begin{definition}\label{radon-measures:definition-narrow}
Finite Radon measures $\mu_k$ \emph{converge narrowly} to a finite Radon measure $\mu$ if $\int f\,d\mu_k \to \int f\,d\mu$ for every $f \in C_b(X)$. A family $\mathcal{M}$ of finite Radon measures is
\emph{tight} if for every $\varepsilon > 0$ there is a compact $K \subseteq X$ with $\mu(X \setminus K) < \varepsilon$ for all $\mu \in \mathcal{M}$.
\end{definition}

Narrow convergence implies weak convergence and, testing with $f = 1$, convergence of the total masses. A single finite Radon measure is tight, by inner regularity; so is a finite family. The
sequences $\delta_k$ and $\frac1k\mathbf{1}_{[0, k]}\lambda$ on $\RR$ are not tight: every compact set lies in some $[-m, m]$, which eventually carries no mass, respectively mass $m/k \to 0$.

\begin{lemma}\label{radon-measures:lemma-weak-to-narrow}
Let $\mu_k \rightharpoonup \mu$ be finite Radon measures on $X$ with $\mu_k(X) \to \mu(X) < \infty$. Then $\mu_k \to \mu$ narrowly.
\end{lemma}

\begin{proof}
Let $f \in C_b(X)$ and $\varepsilon > 0$. By inner regularity choose a compact $K$ with $\mu(X \setminus K) < \varepsilon$, and by Lemma~\ref{general-topology:lemma-lch-shrinking} a function $\phi \in C_c(X)$ with
$0 \leq \phi \leq 1$ and $\phi = 1$ on $K$. Then $\int(1 - \phi)\,d\mu < \varepsilon$, and $\int(1 - \phi)\,d\mu_k = \mu_k(X) - \int\phi\,d\mu_k \to \mu(X) - \int\phi\,d\mu < \varepsilon$, so
$\int(1 - \phi)\,d\mu_k < \varepsilon$ for $k$ large. Since $f\phi \in C_c(X)$, for such $k$
\[\begin{aligned}
\Bigl|\int f\,d\mu_k - \int f\,d\mu\Bigr| &\leq \Bigl|\int f\phi\,d\mu_k - \int f\phi\,d\mu\Bigr| \\
&\quad + \|f\|_\infty\Bigl(\int(1 - \phi)\,d\mu_k + \int(1 - \phi)\,d\mu\Bigr) \\
&\leq \Bigl|\int f\phi\,d\mu_k - \int f\phi\,d\mu\Bigr| + 2\varepsilon\|f\|_\infty,
\end{aligned}\]
and the first term tends to $0$.
\end{proof}

\begin{theorem}[Prokhorov]\label{radon-measures:theorem-prokhorov}
Let $(\mu_k)$ be a tight sequence of Borel probability measures on $X$. Then a subsequence converges narrowly to a Borel probability measure. Conversely, a narrowly convergent sequence of
probability measures on $X$ is tight.
\end{theorem}

\begin{proof}
\emph{Compactness.} The masses on compact sets are bounded by $1$, so by Theorem~\ref{radon-measures:theorem-weak-compactness} a subsequence converges weakly to a Radon measure $\mu$, and
$\mu(X) \leq \liminf_k\mu_k(X) = 1$ by Proposition~\ref{radon-measures:proposition-weak-convergence-sets}(1). Given $\varepsilon > 0$ choose a compact $K$ with $\mu_k(X \setminus K) < \varepsilon$ for all
$k$, and $\phi \in C_c(X)$ with $0 \leq \phi \leq 1$, $\phi = 1$ on $K$. Then $\mu(X) \geq \int\phi\,d\mu = \lim_k\int\phi\,d\mu_k \geq 1 - \varepsilon$. So $\mu(X) = 1$, and the subsequence
converges narrowly by Lemma~\ref{radon-measures:lemma-weak-to-narrow}.

\emph{Tightness.} Let $\mu_k \to \mu$ narrowly and $\varepsilon > 0$. Choose a compact $K$ with $\mu(X \setminus K) < \varepsilon$, and an open $V \supseteq K$ with compact closure
(Lemma~\ref{general-topology:lemma-lch-shrinking}). By Proposition~\ref{radon-measures:proposition-weak-convergence-sets}(1), $\liminf_k\mu_k(V) \geq \mu(V) > 1 - \varepsilon$, so
$\mu_k(X \setminus \overline{V}) < \varepsilon$ for all $k \geq k_0$. The finitely many measures $\mu_1, \ldots, \mu_{k_0}$ are each tight, so a compact set containing $\overline{V}$ and
compact sets for these finitely many measures does the job.
\end{proof}

\begin{example}\label{radon-measures:example-prokhorov}
\begin{enumerate}
\item On $\RR$, a family of probability measures with $\int|x|\,d\mu \leq C$ is tight, by Markov's inequality: $\mu(\RR \setminus [-m, m]) \leq C/m$. So, for instance, the distributions of a sequence of
random variables with bounded means of absolute values have a narrowly convergent subsequence.
\item The Gaussian measures $\gamma_\sigma = (2\pi\sigma^2)^{-1/2}e^{-x^2/2\sigma^2}\lambda$ are tight for $\sigma \leq 1$, and converge narrowly to $\delta_0$ as $\sigma \to 0$; for $\sigma \to \infty$ they are
not tight, and converge weakly to $0$.
\end{enumerate}
\end{example}

Prokhorov proved the theorem for complete separable metric spaces, which need not be locally compact; there every finite Borel measure is tight, and tightness is equivalent to relative compactness of
a family in the topology of narrow convergence \cite[Section 5]{billingsley-convergence}. This is the form used for measures on spaces of paths, where it yields Brownian motion as a limit of random
walks (Part~\ref{part-probability}).

\section{Differentiation of measures}\label{radon-measures:section-differentiation-measures}

The Lebesgue differentiation theorem (Theorem~\ref{lebesgue-differentiation:theorem-lebesgue-differentiation}) says that for $h \in L^1_{\mathrm{loc}}(\RR^n)$ the averages $\frac1{\lambda(B(x, r))}\int_{B(x, r)}h$
tend to $h(x)$ for almost every $x$. In terms of the measure $\nu = h\lambda$ it says that $\nu(B(x, r))/\lambda(B(x, r)) \to h(x)$: the density of $\nu$ can be computed by shrinking balls. We now do this
for an arbitrary Radon measure $\nu$ on $\RR^n$. It has a Lebesgue decomposition $\nu = h\lambda + \nu_s$ with respect to $\lambda$ (Theorem~\ref{lebesgue-differentiation:theorem-lebesgue-decomposition}, which
applies since $\nu$ is $\sigma$-finite by Proposition~\ref{radon-measures:proposition-radon-automatic}), and the shrinking balls recover $h$ and do not see the singular part, except on a null set. In this
section $\lambda$ denotes Lebesgue measure on the Borel sets of $\RR^n$.

\begin{definition}\label{radon-measures:definition-derivative-measure}
Let $\nu$ be a Radon measure on $\RR^n$. The \emph{upper} and \emph{lower derivatives} of $\nu$ with respect to $\lambda$ at $x$ are
\[
\overline{D}\nu(x) = \limsup_{r \to 0}\frac{\nu(B(x, r))}{\lambda(B(x, r))}, \qquad \underline{D}\nu(x) = \liminf_{r \to 0}\frac{\nu(B(x, r))}{\lambda(B(x, r))},
\]
in $[0, \infty]$. If they agree, their common value is the \emph{derivative} $D\nu(x)$.
\end{definition}

\begin{example}\label{radon-measures:example-derivative-measure}
For $\nu = \delta_0$ on $\RR^n$, $\nu(B(x, r)) = 0$ for $r \leq \|x\|$, so $D\nu(x) = 0$ for $x \neq 0$, while $D\nu(0) = \lim_r1/\lambda(B(0, r)) = \infty$. For the measure $\sigma$ on the circle
(Example~\ref{radon-measures:example-positive-functionals}(3)), $D\sigma = 0$ outside $S^1$, because the balls around a point at positive distance from $S^1$ eventually miss $S^1$; at a point of $S^1$ the ball
$B(x, r)$ meets $S^1$ in an arc of length about $2r$, so $\sigma(B(x, r))/\pi r^2$ behaves like $2/(\pi r)$ and $D\sigma = \infty$. In both cases the derivative is $0$ almost everywhere, in accordance
with the singularity of these measures.
\end{example}

\begin{theorem}[Differentiation of measures]\label{radon-measures:theorem-differentiation-measures}
Let $\nu$ be a Radon measure on $\RR^n$, and $\nu = h\lambda + \nu_s$ its Lebesgue decomposition, with $h \geq 0$ Borel and $\nu_s \perp \lambda$. Then $h \in L^1_{\mathrm{loc}}(\RR^n)$, and for
$\lambda$-almost every $x$
\[
\lim_{r \to 0}\frac1{\lambda(B(x, r))}\Bigl(\int_{B(x, r)}|h(y) - h(x)|\,d\lambda(y) + \nu_s(B(x, r))\Bigr) = 0 .
\]
In particular $D\nu = h$ almost everywhere, and $\nu \perp \lambda$ if and only if $D\nu = 0$ almost everywhere.
\end{theorem}

\begin{proof}
For compact $K$, $\int_Kh\,d\lambda \leq \nu(K) < \infty$, so $h \in L^1_{\mathrm{loc}}$, and by the Lebesgue differentiation theorem
(Theorem~\ref{lebesgue-differentiation:theorem-lebesgue-differentiation}) the first term tends to $0$ almost everywhere. It remains to show that $\overline{D}\nu_s = 0$ almost everywhere. Let $S$ be a Borel set with
$\lambda(S) = 0$ and $\nu_s(\RR^n \setminus S) = 0$. The measure $\nu_s$ is a Borel measure finite on compact sets ($\nu_s \leq \nu$), so it is a Radon measure
(Proposition~\ref{radon-measures:proposition-radon-automatic}). Fix $t > 0$ and $R > 0$, and let
\[
A = \{x \in B(0, R) \setminus S : \overline{D}\nu_s(x) > t\}.
\]
We show $\lambda^*(A) = 0$. Let $\varepsilon > 0$. The Borel set $B(0, R) \setminus S$ has $\nu_s$-measure $0$, so by outer regularity it is contained in an open set $U$ with $\nu_s(U) < \varepsilon$. Let $G$ be
the union of all open balls $B$ with $B \subseteq U$ and $\nu_s(B) > t\lambda(B)$. It is open and contains $A$: for $x \in A$ we have $x \in U$, so $B(x, r) \subseteq U$ for all small $r$, and
$\nu_s(B(x, r)) > t\lambda(B(x, r))$ for some of these $r$. Let $K \subseteq G$ be compact. Finitely many of these balls cover $K$, and by the Vitali covering lemma
(Lemma~\ref{lebesgue-differentiation:lemma-vitali-covering}) they contain pairwise disjoint balls $B_1, \ldots, B_k$ with
\[
\lambda(K) \leq 3^n\sum_j\lambda(B_j) < \frac{3^n}{t}\sum_j\nu_s(B_j) \leq \frac{3^n}{t}\nu_s(U) < \frac{3^n\varepsilon}{t}.
\]
By inner regularity of $\lambda$, $\lambda^*(A) \leq \lambda(G) \leq 3^n\varepsilon/t$, and as $\varepsilon$ is arbitrary, $A$ is a null set. Taking the union over $t = 1/j$ and $R = j$, $j \geq 1$, and
adding the null set $S$, we find that $\overline{D}\nu_s = 0$ outside a null set.

Finally, $\nu(B(x, r))/\lambda(B(x, r))$ differs from the average of $h$ over $B(x, r)$ by the quotient $\nu_s(B(x, r))/\lambda(B(x, r))$, and the average differs from $h(x)$ by
at most the average of $|h - h(x)|$; so
$D\nu = h$ almost everywhere. If $\nu \perp \lambda$, then $h = 0$ almost everywhere by uniqueness of the decomposition, so $D\nu = 0$ almost everywhere; conversely, if $D\nu = 0$ almost everywhere,
then $h = 0$ almost everywhere and $\nu = \nu_s$.
\end{proof}

So the Radon--Nikodym derivative of the absolutely continuous part, an object obtained in Theorem~\ref{lebesgue-differentiation:theorem-radon-nikodym} by an abstract Hilbert space argument, can be computed as a
limit of ratios. Balls are not essential: any sets that shrink to $x$ and occupy a fixed proportion of a ball around $x$ will do.

\begin{corollary}\label{radon-measures:corollary-nicely-shrinking}
Let $\nu = h\lambda + \nu_s$ be as in Theorem~\ref{radon-measures:theorem-differentiation-measures} and $c > 0$. For almost every $x$ the following holds: if $E_k$ are Borel sets with
$E_k \subseteq B(x, r_k)$, $r_k \to 0$ and $\lambda(E_k) \geq c\lambda(B(x, r_k))$, then $\nu(E_k)/\lambda(E_k) \to h(x)$.
\end{corollary}

\begin{proof}
Take $x$ as in the theorem. Then
\[\begin{aligned}
\Bigl|\frac{\nu(E_k)}{\lambda(E_k)} - h(x)\Bigr| &\leq \frac1{\lambda(E_k)}\Bigl(\int_{E_k}|h - h(x)|\,d\lambda + \nu_s(E_k)\Bigr) \\
&\leq \frac1{c\lambda(B(x, r_k))}\Bigl(\int_{B(x, r_k)}|h - h(x)|\,d\lambda + \nu_s(B(x, r_k))\Bigr),
\end{aligned}\]
which tends to $0$.
\end{proof}

\begin{corollary}[Lebesgue's theorem on monotone functions]\label{radon-measures:corollary-monotone-differentiable}
Let $F \colon [a, b] \to \RR$ be continuous and nondecreasing. Then $F$ is differentiable almost everywhere, $F^{\prime} \geq 0$ is integrable, and
\[
\int_a^bF^{\prime}\,d\lambda \leq F(b) - F(a),
\]
with equality if and only if $F$ is absolutely continuous.
\end{corollary}

\begin{proof}
Let $\mu_F$ be the measure of Lemma~\ref{lebesgue-differentiation:lemma-stieltjes}, and $\nu(E) = \mu_F(E \cap [a, b])$ for Borel $E \subseteq \RR$; it is a finite Borel measure on $\RR$, hence a Radon measure. Let
$\nu = h\lambda + \nu_s$ be its Lebesgue decomposition, and let $x \in (a, b)$ be a point where Corollary~\ref{radon-measures:corollary-nicely-shrinking} holds with $c = 1/4$. For
$0 < r < \min(b - x, x - a)$, the intervals $(x, x + r]$ and $(x - r, x]$ have length $r = \frac14\lambda(B(x, 2r))$ and lie in $B(x, 2r)$, and
$\nu((x, x + r]) = F(x + r) - F(x)$, $\nu((x - r, x]) = F(x) - F(x - r)$. So both one-sided difference quotients of $F$ at $x$ tend to $h(x)$, and $F^{\prime}(x) = h(x)$. Hence $F^{\prime} = h$ almost
everywhere on $[a, b]$, and
\[
\int_a^bF^{\prime}\,d\lambda = \int_{[a, b]}h\,d\lambda \leq \nu([a, b]) = F(b) - F(a),
\]
with equality if and only if $\nu_s([a, b]) = 0$, that is, $\nu_s = 0$. In that case $F(x) - F(a) = \nu([a, x]) = \int_a^xh\,d\lambda$ for all $x$, so $F$ is absolutely continuous by
Theorem~\ref{lebesgue-differentiation:theorem-absolutely-continuous}. Conversely, if $F$ is absolutely continuous, the same theorem gives $F(b) - F(a) = \int_a^bF^{\prime}$.
\end{proof}

For the Cantor function (Counterexample~\ref{lebesgue-differentiation:counterexample-cantor-function}) the inequality is strict: $\mu_F$ is singular, so $F^{\prime} = 0$ almost everywhere, as we saw directly,
while $F(1) - F(0) = 1$. The corollary holds for all nondecreasing functions, continuous or not: a nondecreasing function has at most countably many jumps, and the jumps contribute a further singular
part to its measure \cite[Chapter 7]{rudin-rca}. Since a function of bounded variation is a difference of two nondecreasing functions, it is differentiable almost everywhere as well.

\begin{remark}\label{radon-measures:remark-besicovitch}
The proof used Lebesgue measure through the Vitali covering lemma, that is, through the fact that $\lambda(3B) = 3^n\lambda(B)$. To differentiate a Radon measure $\nu$ with respect to an arbitrary Radon
measure $\mu$ on $\RR^n$, which may give very different masses to $B$ and $3B$, enlarging balls is not allowed. The \emph{Besicovitch covering theorem} of the next subsection avoids enlargement
altogether, at the price of allowing a bounded overlap; it uses the geometry of $\RR^n$, but no property of any measure. This is the setting of geometric measure theory, where $\mu$ is often a
Hausdorff measure restricted to a set.
\end{remark}

\subsection{The Besicovitch covering theorem}\label{radon-measures:subsection-besicovitch}

Throughout this subsection, for a Radon measure $\mu$ on $\RR^n$ and an arbitrary set $A \subseteq \RR^n$ we write
\[
\mu^*(A) = \inf\{\mu(E) : E \supseteq A \text{ Borel}\} = \inf\{\mu(U) : U \supseteq A \text{ open}\},
\]
the two infima being equal by outer regularity (Proposition~\ref{radon-measures:proposition-radon-automatic}). This outer measure is countably subadditive, and a set with $\mu^*(A) = 0$ lies in a
Borel set of measure $0$. Closed balls are written $\overline{B}(a, r) = \{x : |x - a| \leq r\}$.

\begin{theorem}[Besicovitch covering theorem]\label{radon-measures:theorem-besicovitch}
There is a number $N = N(n)$ with the following property. Let $\mathcal{F}$ be a family of closed balls in $\RR^n$ of positive radius with $\sup\{\operatorname{diam}B : B \in \mathcal{F}\} < \infty$,
and let $A$ be a bounded set such that every point of $A$ is the centre of some ball of $\mathcal{F}$. Then there are countable subfamilies $\mathcal{G}_1, \ldots, \mathcal{G}_N \subseteq \mathcal{F}$, each
consisting of pairwise disjoint balls, with $A \subseteq \bigcup_{i=1}^N\bigcup_{B \in \mathcal{G}_i}B$.
\end{theorem}

\begin{proof}
For each $a \in A$ fix a ball $\overline{B}(a, r(a)) \in \mathcal{F}$; then $R = \sup_ar(a) < \infty$.

\emph{Step 1: greedy choice.} Choose $a_1 \in A$ with $r_1 = r(a_1) \geq \frac34R$, and put $B_j = \overline{B}(a_j, r_j)$. Inductively, let $A_j = A \setminus (B_1 \cup \cdots \cup B_{j-1})$; if $A_j = \emptyset$ stop,
and otherwise choose $a_j \in A_j$ with $r_j = r(a_j) \geq \frac34\sup\{r(a) : a \in A_j\}$. Let $J \leq \infty$ be the number of balls chosen.

\emph{Step 2: properties of the choice.} (a) If $i < j$, then $r_j \leq \frac43r_i$, since $a_j \in A_j \subseteq A_i$ and $r_i \geq \frac34r(a_j)$.
(b) The balls $\overline{B}(a_j, r_j/3)$ are pairwise disjoint: for $i < j$, $a_j \notin B_i$, so $|a_j - a_i| > r_i = \frac{r_i}3 + \frac{2r_i}3 \geq \frac{r_i}3 + \frac{r_j}2 > \frac{r_i}3 + \frac{r_j}3$ by (a).
(c) If $J = \infty$, then $r_j \to 0$: the disjoint balls $\overline{B}(a_j, r_j/3)$ lie in the bounded set $A + \overline{B}(0, R)$, so $\sum_jr_j^n < \infty$.
(d) $A \subseteq \bigcup_jB_j$. This is clear if $J < \infty$. If $J = \infty$ and $a \in A$ lies in no $B_j$, then $a \in A_j$ for all $j$, so $r_j \geq \frac34r(a) > 0$ for all $j$, contradicting (c).

\emph{Step 3: few earlier balls meet a given one.} Fix $k$ and let $I = \{j < k : B_j \cap B_k \neq \emptyset\}$. We show $|I| \leq L = 20^n + (1 + 3\sqrt2)^n$. Split
$I = I_1 \cup I_2$ with $I_1 = \{j \in I : r_j \leq 3r_k\}$ and $I_2 = \{j \in I : r_j > 3r_k\}$.

For $j \in I_1$: $r_j \geq \frac34r_k$ by (a), and $|a_j - a_k| \leq r_j + r_k$, so the ball $\overline{B}(a_j, r_j/3)$, of radius at least $r_k/4$, lies in
$\overline{B}(a_k, r_k + \frac43r_j) \subseteq \overline{B}(a_k, 5r_k)$. These balls are disjoint by (b), so comparing volumes, $|I_1|(r_k/4)^n \leq (5r_k)^n$, and $|I_1| \leq 20^n$.

For $i, j \in I_2$ with $i < j$, let $x = a_i - a_k$, $y = a_j - a_k$ and $\theta$ the angle between them. We claim $\cos\theta < \frac89$. Since $k > j > i$, the point $a_k$ lies outside $B_i$ and $B_j$, and $a_j$
lies outside $B_i$; since $B_k$ meets $B_i$ and $B_j$,
\[
r_i < |x| \leq r_i + r_k, \qquad r_j < |y| \leq r_j + r_k, \qquad |x - y| = |a_j - a_i| > r_i,
\]
and moreover $r_k < \frac13\min(r_i, r_j)$ and $r_j \leq \frac43r_i$. Suppose $\cos\theta \geq \frac89$. If $|y| \geq |x|$, then $0 \leq |y| - |x| < r_j + r_k - r_i \leq \frac13r_i + \frac13r_i$, and
\[
|x - y|^2 = (|y| - |x|)^2 + 2|x||y|(1 - \cos\theta) < \tfrac49r_i^2 + 2\cdot\tfrac43r_i\cdot\tfrac53r_i\cdot\tfrac19 = \tfrac{76}{81}r_i^2 < r_i^2,
\]
using $|x| < \frac43r_i$ and $|y| \leq r_j + r_k < \frac53r_i$; a contradiction. If $|y| < |x|$, then $2|x|\cos\theta - |y| > |x|(2\cos\theta - 1) \geq \frac79|x|$, so
\[\begin{aligned}
|x - y|^2 &= |x|^2 - |y|\bigl(2|x|\cos\theta - |y|\bigr) \leq |x|^2 - 3r_k\cdot\tfrac79|x| \\
&< (r_i + r_k)^2 - \tfrac73r_kr_i = r_i^2 - \tfrac13r_kr_i + r_k^2 < r_i^2,
\end{aligned}\]
using $|y| > r_j > 3r_k$, $r_i < |x| \leq r_i + r_k$ and $r_k < \frac13r_i$; again a contradiction. So the unit vectors $u_j = (a_j - a_k)/|a_j - a_k|$, $j \in I_2$, satisfy
$|u_i - u_j|^2 = 2 - 2\cos\theta > \frac29$. The balls of radius $\frac{\sqrt2}6$ around them are disjoint and lie in $\overline{B}(0, 1 + \frac{\sqrt2}6)$, so comparing volumes,
$|I_2| \leq \bigl((1 + \frac{\sqrt2}6)/\frac{\sqrt2}6\bigr)^n = (1 + 3\sqrt2)^n$.

\emph{Step 4: colouring.} Let $N = L + 1$. Distribute the balls $B_1, B_2, \ldots$ in turn among $N$ families: put $B_k$ into a family containing no $B_j$, $j < k$, that meets $B_k$. This is possible,
since by Step 3 at most $L$ earlier balls meet $B_k$. Each family consists of pairwise disjoint balls, and by (d) their union covers $A$.
\end{proof}

\begin{counterexample}\label{radon-measures:counterexample-besicovitch}
\begin{enumerate}
\item The balls must be centred at the points to be covered. In $\RR$, let $A = (0, 1]$ and let $\mathcal{F}$ consist of the intervals $[1/k, 2]$, $k \geq 1$, whose centres lie in $(1, \frac32]$. They
cover $A$, but they all contain $1$, so a disjoint subfamily has at most one member, and $N$ disjoint subfamilies cover only $[1/k, 2]$ for some $k$, never $A$.
\item Some overlap cannot be avoided. Let $A = \{-1, 1\}$ and $\mathcal{F} = \{[-2, 0], [0, 2]\}$, the balls of radius $1$ centred at the two points. Covering $A$ needs both intervals, and they meet
at $0$; so one disjoint family does not suffice.
\item Some bound is needed when $A$ is unbounded. In $\RR$ let $A = \ZZ_{\geq 0}$ and $\mathcal{F} = \{\overline{B}(k, k) : k \geq 1\} \cup \{\overline{B}(0, 1)\}$. Every ball contains $0$, so a
disjoint family has at most one member, and finitely many disjoint families cover a bounded set, never $A$.
\end{enumerate}
\end{counterexample}

The constant of the proof is far from optimal; the least possible $N(n)$, the \emph{Besicovitch constant}, is known only in low dimensions. The theorem uses the Euclidean structure (Step 3
measures angles), and it fails in some metric spaces; there the doubling condition of the next subsection is the usual substitute. Besicovitch's theorem gives Vitali-type coverings for every
Radon measure.

\begin{corollary}\label{radon-measures:corollary-besicovitch-vitali}
Let $\mu$ be a Radon measure on $\RR^n$, $A \subseteq \RR^n$ a bounded set, and $\mathcal{F}$ a family of closed balls of positive radius such that every $a \in A$ is the centre of balls of
$\mathcal{F}$ of arbitrarily small radius. Then for every open $U \supseteq A$ there is a countable family of pairwise disjoint balls $B_i \in \mathcal{F}$ contained in $U$ with
$\mu^*(A \setminus \bigcup_iB_i) = 0$.
\end{corollary}

\begin{proof}
Replacing $U$ by its intersection with a large ball, we may assume $U$ bounded, and we may assume $\mu^*(A) > 0$. Let $N$ be as in the theorem and $\theta = \frac1{4N}$.

\emph{Step 1: a fixed proportion.} If $A^{\prime} \subseteq A$ has $\mu^*(A^{\prime}) > 0$ and $U^{\prime} \supseteq A^{\prime}$ is a bounded open set, there are finitely many disjoint balls of $\mathcal{F}$ in
$U^{\prime}$ with $\mu^*(A^{\prime} \setminus \bigcup B_i) \leq (1 - \theta)\mu^*(A^{\prime})$. Indeed, replacing $U^{\prime}$ by a smaller open set we may assume
$\mu(U^{\prime}) \leq (1 + \theta)\mu^*(A^{\prime})$. Apply Theorem~\ref{radon-measures:theorem-besicovitch} to the balls of $\mathcal{F}$ centred in $A^{\prime}$, contained in $U^{\prime}$, of radius at most
$1$: we get $N$ disjoint families covering $A^{\prime}$, so one of them, $\mathcal{G}$, has $\mu(\bigcup_{\mathcal{G}}B) \geq \mu^*(A^{\prime})/N$, and a finite subfamily $B_1, \ldots, B_m$ has
$\mu(\bigcup B_i) \geq \mu^*(A^{\prime})/(2N)$. Then
\[\begin{aligned}
\mu^*\Bigl(A^{\prime} \setminus \bigcup_iB_i\Bigr) &\leq \mu\Bigl(U^{\prime} \setminus \bigcup_iB_i\Bigr) = \mu(U^{\prime}) - \mu\Bigl(\bigcup_iB_i\Bigr) \\
&\leq \Bigl(1 + \theta - \frac1{2N}\Bigr)\mu^*(A^{\prime}) = (1 - \theta)\mu^*(A^{\prime}).
\end{aligned}\]

\emph{Step 2: iteration.} Apply Step 1 to $A$ and $U$, then to the remainder $A^{\prime} = A \setminus (B_1 \cup \cdots \cup B_m)$ and the open set $U^{\prime} = U \setminus (B_1 \cup \cdots \cup B_m)$, and so on,
stopping if the remainder becomes $\mu^*$-null. The balls chosen are pairwise disjoint and lie in $U$, and after $k$ steps the uncovered part of $A$ has outer measure at most
$(1 - \theta)^k\mu^*(A) \to 0$.
\end{proof}

\begin{definition}\label{radon-measures:definition-derivative-radon}
Let $\mu$ and $\nu$ be Radon measures on $\RR^n$, and $S = \operatorname{supp}\mu$, so that $\mu(\overline{B}(x, r)) > 0$ for $x \in S$ and $r > 0$, and $\mu(\RR^n \setminus S) = 0$. For $x \in S$ put
\[
\overline{D}_\mu\nu(x) = \limsup_{r \to 0}\frac{\nu(\overline{B}(x, r))}{\mu(\overline{B}(x, r))}, \qquad \underline{D}_\mu\nu(x) = \liminf_{r \to 0}\frac{\nu(\overline{B}(x, r))}{\mu(\overline{B}(x, r))},
\]
and, where they agree, write $D_\mu\nu(x)$ for the common value, the \emph{derivative of $\nu$ with respect to $\mu$}.
\end{definition}

\begin{theorem}[Differentiation of Radon measures]\label{radon-measures:theorem-besicovitch-differentiation}
Let $\mu$ and $\nu$ be Radon measures on $\RR^n$. Then $D_\mu\nu(x)$ exists and is finite for $\mu$-almost every $x$, and it is a Borel function (defined $\mu$-almost everywhere). If
$\nu = \nu_a + \nu_s$ is the Lebesgue decomposition with respect to $\mu$, then $\nu_a = (D_\mu\nu)\mu$ and $D_\mu\nu_s = 0$ $\mu$-almost everywhere. In particular $\nu \ll \mu$ if and only if
$\nu(E) = \int_ED_\mu\nu\,d\mu$ for all Borel sets $E$.
\end{theorem}

\begin{proof}
\emph{Step 1: two comparison inequalities.} Let $0 < \alpha < \infty$ and $A \subseteq S$ bounded. (i) If $\underline{D}_\mu\nu(x) \leq \alpha$ for all $x \in A$, then $\nu^*(A) \leq \alpha\mu^*(A)$. (ii) If
$\overline{D}_\mu\nu(x) \geq \alpha$ for all $x \in A$, then $\nu^*(A) \geq \alpha\mu^*(A)$. For (i), let $\varepsilon > 0$ and $U \supseteq A$ open with $\mu(U) \leq \mu^*(A) + \varepsilon$. The closed balls
$\overline{B}(a, r)$ with $a \in A$ and $\nu(\overline{B}(a, r)) \leq (\alpha + \varepsilon)\mu(\overline{B}(a, r))$ have arbitrarily small radii at every $a \in A$, so by
Corollary~\ref{radon-measures:corollary-besicovitch-vitali}, applied to the measure $\nu$, disjoint ones $B_i \subseteq U$ cover $\nu$-almost all of $A$. Then
\[
\nu^*(A) \leq \sum_i\nu(B_i) \leq (\alpha + \varepsilon)\sum_i\mu(B_i) \leq (\alpha + \varepsilon)\mu(U) \leq (\alpha + \varepsilon)(\mu^*(A) + \varepsilon),
\]
and $\varepsilon \to 0$ gives (i). For (ii), let $\varepsilon < \alpha$, take $U \supseteq A$ open with $\nu(U) \leq \nu^*(A) + \varepsilon$, and balls with $\nu(B) \geq (\alpha - \varepsilon)\mu(B)$; by the
corollary, applied to $\mu$, disjoint ones $B_i \subseteq U$ cover $\mu$-almost all of $A$, and $(\alpha - \varepsilon)\mu^*(A) \leq (\alpha - \varepsilon)\sum_i\mu(B_i) \leq \nu(U) \leq \nu^*(A) + \varepsilon$.

\emph{Step 2: the limit exists and is finite.} For rationals $0 < a < b$ and $m \geq 1$, the set $A = \{x \in S : \underline{D}_\mu\nu(x) < a < b < \overline{D}_\mu\nu(x)\} \cap B(0, m)$ satisfies
$b\mu^*(A) \leq \nu^*(A) \leq a\mu^*(A)$ by Step 1, hence $\mu^*(A) = 0$, as $\mu^*(A) < \infty$. Similarly $A = \{\overline{D}_\mu\nu = \infty\} \cap B(0, m)$ satisfies
$\alpha\mu^*(A) \leq \nu(B(0, m)) < \infty$ for all $\alpha$, so $\mu^*(A) = 0$. Taking countable unions, $\underline{D}_\mu\nu = \overline{D}_\mu\nu < \infty$ $\mu$-almost everywhere.

\emph{Step 3: measurability.} For fixed $x$, the functions $r \mapsto \mu(\overline{B}(x, r))$ and $r \mapsto \nu(\overline{B}(x, r))$ are right-continuous, because $\overline{B}(x, r)$ is the decreasing
intersection of the $\overline{B}(x, s)$, $s > r$; so $\overline{D}_\mu\nu$ and $\underline{D}_\mu\nu$ can be computed along rational $r$. For fixed $r$, $x \mapsto \mu(\overline{B}(x, r))$ is upper
semicontinuous: if $x_k \to x$, then $\limsup_k\mathbf{1}_{\overline{B}(x_k, r)} \leq \mathbf{1}_{\overline{B}(x, r)}$ pointwise, and Fatou's lemma applied to $\mathbf{1}_K - \mathbf{1}_{\overline{B}(x_k, r)}$,
for a compact $K$ containing all these balls, gives $\limsup_k\mu(\overline{B}(x_k, r)) \leq \mu(\overline{B}(x, r))$. So these functions are Borel, and so are $\overline{D}_\mu\nu$ and $\underline{D}_\mu\nu$
($\limsup$ and $\liminf$ of countably many Borel functions on the closed set $S$) and the set where they agree and are finite.

\emph{Step 4: singular measures.} If $\nu \perp \mu$, let $Z$ be a Borel set with $\mu(Z) = 0 = \nu(\RR^n \setminus Z)$. For $\alpha > 0$ and $m \geq 1$, the set
$A = \{\overline{D}_\mu\nu \geq \alpha\} \cap B(0, m) \setminus Z$ has $\alpha\mu^*(A) \leq \nu^*(A) = 0$ by (ii). So $D_\mu\nu = 0$ $\mu$-almost everywhere.

\emph{Step 5: absolutely continuous measures.} Let $\nu \ll \mu$ and $E$ Borel. The set $\{D_\mu\nu = 0\} \cap B(0, m)$ has $\nu$-measure at most $\alpha\mu(B(0, m))$ for every $\alpha > 0$ by (i), so
$\nu(\{D_\mu\nu = 0\}) = 0$; the $\mu$-null set where $D_\mu\nu$ is undefined or infinite is $\nu$-null too. Fix $t > 1$ and let $E_m = E \cap \{t^m \leq D_\mu\nu < t^{m+1}\}$, $m \in \ZZ$, which
cover $E$ up to a $\nu$-null set on which $D_\mu\nu$ is $0$ or undefined. By Step 1, applied to the bounded pieces $E_m \cap B(0, j)$ and letting $j \to \infty$,
\[
\nu(E_m) \leq t^{m+1}\mu(E_m) \leq t\int_{E_m}D_\mu\nu\,d\mu, \qquad \nu(E_m) \geq t^m\mu(E_m) \geq t^{-1}\int_{E_m}D_\mu\nu\,d\mu .
\]
Summing over $m$, $t^{-1}\int_ED_\mu\nu\,d\mu \leq \nu(E) \leq t\int_ED_\mu\nu\,d\mu$, and letting $t \to 1$, $\nu(E) = \int_ED_\mu\nu\,d\mu$.

\emph{Step 6: general $\nu$.} Radon measures on $\RR^n$ are $\sigma$-finite, so the Lebesgue decomposition $\nu = \nu_a + \nu_s$ exists (Theorem~\ref{signed-measures:theorem-lebesgue-radon-nikodym}); $\nu_a$ and
$\nu_s$ are Borel measures with $\nu_a, \nu_s \leq \nu$, hence Radon (Proposition~\ref{radon-measures:proposition-radon-automatic}). The quotients are additive in $\nu$, so
$D_\mu\nu = D_\mu\nu_a + D_\mu\nu_s = D_\mu\nu_a$ $\mu$-almost everywhere by Step 4, and $\nu_a = (D_\mu\nu_a)\mu = (D_\mu\nu)\mu$ by Step 5. The last assertion follows: if $\nu \ll \mu$, then $\nu = \nu_a$;
conversely, $\nu = (D_\mu\nu)\mu$ is absolutely continuous with respect to $\mu$.
\end{proof}

\begin{example}\label{radon-measures:example-besicovitch-differentiation}
\begin{enumerate}
\item For $\mu = \lambda$ the theorem recovers Theorem~\ref{radon-measures:theorem-differentiation-measures}, with closed instead of open balls, by a proof that does not use the Lebesgue differentiation
theorem.
\item Let $\mu = \delta_0 + \lambda$ on $\RR$ and $\nu = \lambda$. Then $\nu \ll \mu$, with density $\mathbf{1}_{\RR \setminus \{0\}}$. At $x \neq 0$ the quotient is $2r/(2r) = 1$ for small $r$, and at $0$ it
is $2r/(1 + 2r) \to 0$: the derivative detects the atom, a $\mu$-non-null point where the density vanishes.
\item Let $\mu = \mu_F$ for the Cantor function $F$, restricted to $[0, 1]$, and $\nu = \lambda$ on $[0, 1]$. We saw after Corollary~\ref{radon-measures:corollary-monotone-differentiable} that $\mu$ is singular,
and indeed $D_\mu\lambda = 0$ $\mu$-almost everywhere: a point $x$ of the Cantor set lies in one of the $2^{k+1}$ intervals of length $3^{-k-1}$ of the $(k+1)$-st stage, which has $\mu$-mass $2^{-k-1}$ and lies in
$\overline{B}(x, r)$ for $r \geq 3^{-k-1}$; so for $3^{-k-1} \leq r \leq 3^{-k}$ the quotient $\lambda(\overline{B}(x, r))/\mu(\overline{B}(x, r))$ is at most $2\cdot3^{-k}/2^{-k-1} = 4(2/3)^k$, which tends to $0$.
\end{enumerate}
\end{example}

The theorem is the measure-theoretic form of the Radon--Nikodym theorem: the derivative $d\nu_a/d\mu$, obtained abstractly in Chapter~\ref{signed-measures}, is a limit of quotients of masses of small
balls. Applied to $\nu = f^{\pm}\mu$ it shows that for $f \in L^1_{\mathrm{loc}}(\mu)$ the $\mu$-averages of $f$ over $\overline{B}(x, r)$ tend to $f(x)$ for $\mu$-almost every $x$; so the Lebesgue
differentiation theorem holds for every Radon measure on $\RR^n$ \cite[Section 1.7]{evans-gariepy}.


\subsection{Differentiation with respect to doubling measures}\label{radon-measures:subsection-doubling}

Densities measure a measure $\mu$ against the function $r \mapsto \alpha(s)(2r)^s$. Measuring it against another measure $\nu$ instead, through the ratios $\mu(\overline{B}(x, r))/\nu(\overline{B}(x, r))$,
leads to the general form of the Radon--Nikodym theorem in which the derivative $d\mu/d\nu$ is computed as a limit, as in Theorem~\ref{radon-measures:theorem-differentiation-measures} for
$\nu = \lambda$ and in Theorem~\ref{radon-measures:theorem-besicovitch-differentiation} for every Radon measure on $\RR^n$, whose proof used the Euclidean geometry through Besicovitch's theorem. The proof there used $\lambda(3B) = 3^n\lambda(B)$, and the $5r$-covering lemma (Lemma~\ref{lebesgue-differentiation:lemma-five-r-covering}) shows that a weak form of this property is all one
needs, in any metric space. Throughout, $X$ is a nonempty metric space, and for a Borel measure $\nu$ on $X$ and any $A \subseteq X$ we write $\nu^*(A) = \inf\{\nu(E) : E \supseteq A \text{ Borel}\}$.
This outer measure is countably subadditive (take a union of Borel supersets), and a set with $\nu^*(A) = 0$ lies in a Borel set of measure $0$.

\begin{definition}\label{radon-measures:definition-doubling}
A Borel measure $\nu$ on $X$ is \emph{doubling} if $0 < \nu(\overline{B}(x, r)) < \infty$ for all $x \in X$ and $r > 0$, and there is a constant $C$ with
\[
\nu(\overline{B}(x, 2r)) \leq C\,\nu(\overline{B}(x, r)) \qquad \text{for all } x \in X,\ r > 0 .
\]
\end{definition}

Iterating, $\nu(\overline{B}(x, 2^kr)) \leq C^k\nu(\overline{B}(x, r))$; in particular, for the enlarged ball $\widehat{B} = \overline{B}(x, 5r)$ of $B = \overline{B}(x, r)$,
\begin{equation}\label{radon-measures:equation-doubling-five}
\nu(\widehat{B}) \leq \nu(\overline{B}(x, 8r)) \leq C^3\nu(B) .
\end{equation}
A doubling measure is finite on bounded sets, since every bounded set lies in a closed ball, and hence $\sigma$-finite: $X = \bigcup_k\overline{B}(x_0, k)$ for any $x_0 \in X$.

\begin{example}\label{radon-measures:example-doubling}
\begin{enumerate}
\item Lebesgue measure on $\RR^n$ is doubling with $C = 2^n$, since $\lambda(\overline{B}(x, r)) = \alpha(n)(2r)^n$.
\item If $w$ is a Borel function on $\RR^n$ with $0 < m \leq w \leq M$, then $\nu = w\lambda$ is doubling with $C = 2^nM/m$:
$\nu(\overline{B}(x, 2r)) \leq M\lambda(\overline{B}(x, 2r)) = 2^nM\lambda(\overline{B}(x, r)) \leq 2^n(M/m)\nu(\overline{B}(x, r))$.
\item The Gaussian measure $\gamma = e^{-y^2/2}\lambda$ on $\RR$ is not doubling. For $x > 2$ the density $e^{-y^2/2}$ is decreasing on $[x - 2, x + 2]$, so
\[
\frac{\gamma([x - 2, x + 2])}{\gamma([x - 1, x + 1])} \geq \frac{\frac12e^{-(x - 3/2)^2/2}}{2e^{-(x - 1)^2/2}} = \frac14e^{(x - 5/4)/2},
\]
estimating the numerator by the integral over $[x - 2, x - \frac32]$; the right side is unbounded as $x \to \infty$. Far out, the measure of a ball is concentrated near its left end.
\item A point mass $\delta_a$ on $\RR$ is not doubling, since balls missing $a$ have measure $0$; but $\delta_a$ is doubling on the one-point metric space $\{a\}$. Doubling is a property of the measure
together with the metric space it lives on.
\end{enumerate}
\end{example}

We need one approximation lemma, which replaces the role of $C_c(\RR^n)$ in the Lebesgue differentiation theorem.

\begin{lemma}\label{radon-measures:lemma-lipschitz-dense}
Let $\nu$ be a finite Borel measure on a metric space $X$ and $1 \leq p < \infty$. The bounded Lipschitz functions are dense in $L^p(\nu)$.
\end{lemma}

\begin{proof}
By Proposition~\ref{measure-theory:proposition-measurable-functions}(4) and dominated convergence, simple functions $\sum_ia_i\mathbf{1}_{E_i}$ with $E_i$ Borel are dense, so it suffices to approximate
$\mathbf{1}_E$ for a Borel set $E$. Let $\varepsilon > 0$. By Lemma~\ref{radon-measures:lemma-closed-approximation} there is a closed set $F \subseteq E$ with $\nu(E \setminus F) < \varepsilon$, so
$\|\mathbf{1}_E - \mathbf{1}_F\|_p^p < \varepsilon$. If $F = \emptyset$, then $\|\mathbf{1}_E\|_p^p < \varepsilon$ and $0$ approximates. Otherwise let $g_k(x) = \max(0, 1 - k\,d(x, F))$. It is Lipschitz with constant $k$, takes values in $[0, 1]$,
equals $1$ on $F$, and $g_k(x) \to 0$ for $x \notin F$ because $d(x, F) > 0$ for such $x$, $F$ being closed. So $g_k \to \mathbf{1}_F$ pointwise with $|g_k - \mathbf{1}_F| \leq 1$, and as $\nu$ is finite,
dominated convergence gives $\|g_k - \mathbf{1}_F\|_p \to 0$.
\end{proof}

From now on $\nu$ is a doubling measure on $X$ with constant $C$.

\begin{theorem}[Maximal inequality for doubling measures]\label{radon-measures:theorem-doubling-maximal}
Let $f \in L^1(\nu)$ and $t > 0$, and let $A$ be the set of $x \in X$ such that
\[
\frac{1}{\nu(\overline{B}(x, r))}\int_{\overline{B}(x, r)}|f|\,d\nu > t \qquad \text{for some } 0 < r \leq 1 .
\]
Then $\nu^*(A) \leq C^3\|f\|_1/t$.
\end{theorem}

\begin{proof}
Let $\mathcal{F}$ be the family of balls $B = \overline{B}(x, r)$ with $x \in A$, $0 < r \leq 1$ and $\int_B|f|\,d\nu > t\nu(B)$; every point of $A$ is the centre of a ball of $\mathcal{F}$. Let
$\mathcal{G}$ be the disjoint subfamily of Lemma~\ref{lebesgue-differentiation:lemma-five-r-covering}. For finitely many balls of $\mathcal{G}$ the integrals $\int_B|f|\,d\nu$ add up to at most $\|f\|_1$, and each of them is
$> t\nu(B) > 0$; so for each $m$ only finitely many exceed $1/m$, and $\mathcal{G} = \{B_1, B_2, \ldots\}$ is countable. Every ball of $\mathcal{F}$ lies in some $\widehat{B_i}$, so
$A \subseteq \bigcup_i\widehat{B_i}$, and by \eqref{radon-measures:equation-doubling-five}
\[
\nu^*(A) \leq \sum_i\nu(\widehat{B_i}) \leq C^3\sum_i\nu(B_i) < \frac{C^3}{t}\sum_i\int_{B_i}|f|\,d\nu \leq \frac{C^3}{t}\|f\|_1 .
\]
(If $A$ is empty there is nothing to prove.)
\end{proof}

\begin{theorem}[Lebesgue differentiation theorem for doubling measures]\label{radon-measures:theorem-doubling-lebesgue}
Let $f \colon X \to \RR$ be a Borel function with $\int_B|f|\,d\nu < \infty$ for every closed ball $B$. For $\nu$-almost every $x \in X$,
\[
\lim_{r \to 0}\frac{1}{\nu(\overline{B}(x, r))}\int_{\overline{B}(x, r)}|f(y) - f(x)|\,d\nu(y) = 0,
\]
and in particular the averages of $f$ over $\overline{B}(x, r)$ tend to $f(x)$.
\end{theorem}

\begin{proof}
Write $\langle g \rangle_{x, r}$ for the average of a function $g$ over $\overline{B}(x, r)$ with respect to $\nu$, and $\Omega(x) = \limsup_{r \to 0}\langle|f - f(x)|\rangle_{x, r}$. Fix $x_0 \in X$, an integer $k \geq 1$,
$t > 0$ and $\varepsilon > 0$, and let $D = \overline{B}(x_0, k + 1)$. The measure $E \mapsto \nu(E \cap D)$ is finite, so by Lemma~\ref{radon-measures:lemma-lipschitz-dense} there is a bounded continuous $g$ with
$\int_D|f - g|\,d\nu < \varepsilon$. Let $h = (f - g)\mathbf{1}_D$, so $\|h\|_1 < \varepsilon$.

Let $x \in B(x_0, k)$ and $0 < r \leq 1$. Then $\overline{B}(x, r) \subseteq D$, and for $y \in \overline{B}(x, r)$
\[
|f(y) - f(x)| \leq |h(y)| + |h(x)| + |g(y) - g(x)| .
\]
Averaging over $y$ gives $\langle|f - f(x)|\rangle_{x, r} \leq \langle|h|\rangle_{x, r} + |h(x)| + \sup_{\overline{B}(x, r)}|g - g(x)|$, and the last term tends to $0$ as $r \to 0$ because $g$ is continuous. So, if
$\Omega(x) > 2t$, then either $\langle|h|\rangle_{x, r} > t$ for some $r \leq 1$, or $|h(x)| > t$. By Theorem~\ref{radon-measures:theorem-doubling-maximal} and Markov's inequality
(Corollary~\ref{measure-theory:corollary-integral-properties}(4)),
\[
\nu^*\bigl(\{x \in B(x_0, k) : \Omega(x) > 2t\}\bigr) \leq \frac{C^3\varepsilon}{t} + \frac{\varepsilon}{t} .
\]
As $\varepsilon$ is arbitrary, this set has outer measure $0$. The set where $\Omega > 0$ is the union of these sets over $k \geq 1$ and $t = 1/j$, $j \geq 1$, so it has outer measure $0$ and lies in a
$\nu$-null Borel set. Finally $|\langle f \rangle_{x, r} - f(x)| \leq \langle|f - f(x)|\rangle_{x, r}$.
\end{proof}

\begin{theorem}[Differentiation of measures with respect to a doubling measure]\label{radon-measures:theorem-doubling-differentiation}
Let $\mu$ be a Borel measure on $X$ that is finite on bounded Borel sets, and let $\mu = h\nu + \mu_s$ be its Lebesgue decomposition with respect to $\nu$
(Theorem~\ref{lebesgue-differentiation:theorem-lebesgue-decomposition}), with $h \geq 0$ Borel and $\mu_s \perp \nu$. Then for $\nu$-almost every $x$
\[
\lim_{r \to 0}\frac{\mu(\overline{B}(x, r))}{\nu(\overline{B}(x, r))} = h(x) .
\]
In particular, if $\mu \ll \nu$, the Radon--Nikodym derivative is $\frac{d\mu}{d\nu}(x) = \lim_{r \to 0}\mu(\overline{B}(x, r))/\nu(\overline{B}(x, r))$ for $\nu$-almost every $x$; and $\mu \perp \nu$ if and only
if this limit is $0$ $\nu$-almost everywhere.
\end{theorem}

\begin{proof}
Both measures are $\sigma$-finite, $X$ being the union of the balls $\overline{B}(x_0, k)$, so the decomposition exists. For every closed ball $B$, $\int_Bh\,d\nu \leq \mu(B) < \infty$, so by
Theorem~\ref{radon-measures:theorem-doubling-lebesgue} the averages of $h$ over $\overline{B}(x, r)$ tend to $h(x)$ for almost every $x$. Since
\[
\frac{\mu(\overline{B}(x, r))}{\nu(\overline{B}(x, r))} = \frac{1}{\nu(\overline{B}(x, r))}\int_{\overline{B}(x, r)}h\,d\nu + \frac{\mu_s(\overline{B}(x, r))}{\nu(\overline{B}(x, r))},
\]
it remains to show that the last ratio tends to $0$ for almost every $x$.

Let $S$ be a Borel set with $\nu(S) = 0$ and $\mu_s(X \setminus S) = 0$. Fix $x_0 \in X$, $k \geq 1$, $t > 0$ and $\varepsilon > 0$, and let $A$ be the set of $x \in B(x_0, k) \setminus S$ with
$\limsup_{r \to 0}\mu_s(\overline{B}(x, r))/\nu(\overline{B}(x, r)) > t$. The measure $E \mapsto \mu_s(E \cap B(x_0, k + 1))$ is finite and vanishes on the Borel set $B(x_0, k) \setminus S$, so by
Lemma~\ref{radon-measures:lemma-closed-approximation} there is an open $U^{\prime} \supseteq B(x_0, k) \setminus S$ with $\mu_s(U^{\prime} \cap B(x_0, k + 1)) < \varepsilon$. The set $U = U^{\prime} \cap B(x_0, k + 1)$ is
open, contains $A$, and has $\mu_s(U) < \varepsilon$.

Let $\mathcal{F}$ be the family of balls $B = \overline{B}(x, r)$ with $x \in A$, $r \leq 1$, $B \subseteq U$ and $\mu_s(B) > t\nu(B)$. Each $x \in A$ is the centre of such balls, because
$\overline{B}(x, r) \subseteq U$ for small $r$. Let $\mathcal{G}$ be the disjoint subfamily of Lemma~\ref{lebesgue-differentiation:lemma-five-r-covering}. Its balls lie in the bounded set $U$, which has finite
$\nu$-measure, and have positive $\nu$-measure, so $\mathcal{G} = \{B_1, B_2, \ldots\}$ is countable, as in the proof of Theorem~\ref{radon-measures:theorem-doubling-maximal}. As there,
$A \subseteq \bigcup_i\widehat{B_i}$ and
\[
\nu^*(A) \leq C^3\sum_i\nu(B_i) < \frac{C^3}{t}\sum_i\mu_s(B_i) \leq \frac{C^3}{t}\mu_s(U) < \frac{C^3\varepsilon}{t} .
\]
So $\nu^*(A) = 0$. Taking the union over $k$ and $t = 1/j$, and adding the null set $S$, the ratio $\mu_s(\overline{B}(x, r))/\nu(\overline{B}(x, r))$ tends to $0$ outside a $\nu$-null set.

If $\mu \ll \nu$, then $\mu_s = 0$ by uniqueness of the decomposition, and $h = d\mu/d\nu$ is the Radon--Nikodym derivative (Theorem~\ref{lebesgue-differentiation:theorem-radon-nikodym}). If $\mu \perp \nu$, then
$h\nu = 0$ by uniqueness, so $h = 0$ almost everywhere; conversely, if the limit is $0$ almost everywhere, then $h = 0$ almost everywhere and $\mu = \mu_s$.
\end{proof}

So the Radon--Nikodym derivative, obtained abstractly in Theorem~\ref{lebesgue-differentiation:theorem-radon-nikodym}, is a limit of ratios of measures of balls whenever the reference measure is doubling. For
$f \in L^1_{\mathrm{loc}}$ and $\mu = f\nu$ the theorem is the Lebesgue differentiation theorem again; for $\mu = \nu\llcorner E$, $E$ Borel, it is a density theorem:
$\nu(E \cap \overline{B}(x, r))/\nu(\overline{B}(x, r))$ tends to $1$ at almost every point of $E$ and to $0$ at almost every point outside $E$.

\begin{remark}\label{radon-measures:remark-doubling}
Doubling is sufficient, not necessary. On $\RR^n$ every Radon measure can serve as $\nu$, by the Besicovitch covering theorem
(Remark~\ref{radon-measures:remark-besicovitch}); so the non-doubling Gaussian measure of Example~\ref{radon-measures:example-doubling}(3)
still differentiates. In general metric spaces some covering property is needed; spaces with a doubling measure, the \emph{spaces of homogeneous type}, are the standard setting for analysis on metric
spaces \cite[Chapter 1]{heinonen-2001}.
\end{remark}

\begin{exercise}\label{radon-measures:exercise-doubling-totally-bounded}
Let $\nu$ be a doubling measure on $X$ with constant $C$.
\begin{enumerate}
\item Let $0 < \varepsilon \leq R$ and let $m \geq 0$ be an integer with $2^m\varepsilon \geq 2R$. Show that a set of points of $\overline{B}(x, R)$ with pairwise distances $> 2\varepsilon$ has at most $C^{m+1}$
elements.
\item Deduce that closed balls in $X$ are totally bounded, that $X$ is separable, and that if $X$ is complete, closed balls are compact.
\item Deduce that the Hilbert space $\ell^2$ of square-summable real sequences carries no doubling measure.
\end{enumerate}
\end{exercise}

\begin{solution}
(1) Let $y_1, \ldots, y_N \in \overline{B}(x, R)$ with $d(y_i, y_j) > 2\varepsilon$ for $i \neq j$. The balls $\overline{B}(y_i, \varepsilon)$ are pairwise disjoint and lie in $\overline{B}(x, R + \varepsilon) \subseteq \overline{B}(x, 2R)$.
Moreover $\overline{B}(x, R) \subseteq \overline{B}(y_i, 2R)$, so
$\nu(\overline{B}(x, R)) \leq \nu(\overline{B}(y_i, 2R)) \leq C^m\nu(\overline{B}(y_i, 2^{-m}2R)) \leq C^m\nu(\overline{B}(y_i, \varepsilon))$. Summing over $i$,
\[
NC^{-m}\nu(\overline{B}(x, R)) \leq \sum_i\nu(\overline{B}(y_i, \varepsilon)) \leq \nu(\overline{B}(x, 2R)) \leq C\nu(\overline{B}(x, R)),
\]
and dividing by $\nu(\overline{B}(x, R)) > 0$ gives $N \leq C^{m+1}$.

(2) Given $\varepsilon \leq R$, choose a set $\{y_1, \ldots, y_N\} \subseteq \overline{B}(x, R)$ with pairwise distances $> 2\varepsilon$ and with $N$ as large as possible, which is possible by (1). Every
$y \in \overline{B}(x, R)$ has $d(y, y_i) \leq 2\varepsilon$ for some $i$, since otherwise $y$ could be added. So finitely many open balls of radius $3\varepsilon$ cover $\overline{B}(x, R)$; for
$\varepsilon > R$ the single ball $B(x, 3\varepsilon)$ does. Hence $\overline{B}(x, R)$ is totally bounded. A totally bounded metric space is separable (the centres of finite $1/j$-nets, $j \geq 1$, form a
dense countable set), and $X = \bigcup_k\overline{B}(x_0, k)$ is a countable union of such spaces, so $X$ is separable. A closed subset of a complete space is complete, so if $X$ is complete, closed balls
are compact by Theorem~\ref{general-topology:theorem-compact-metric}.

(3) The unit vectors $e_1, e_2, \ldots$ lie in the closed unit ball of $\ell^2$ and have pairwise distances $\|e_i - e_j\| = \sqrt2$. So the closed unit ball is not totally bounded: no finitely many balls of
radius $\frac12$ cover it, since each such ball contains at most one $e_i$. By (2), there is no doubling measure on $\ell^2$. Doubling measures live on spaces that are ``finite-dimensional'' in a metric sense.
\end{solution}

\section{Exercises}\label{radon-measures:section-exercises}

\begin{exercise}\label{radon-measures:exercise-riesz-circle}
Let $\sigma$ be the measure on $\RR^2$ of Example~\ref{radon-measures:example-positive-functionals}(3), and $\varphi(t) = (\cos t, \sin t)$. Show that $\sigma(E) = \lambda(\{t \in [0, 2\pi] : \varphi(t) \in E\})$
for every Borel set $E$, and deduce that an arc $\{\varphi(t) : \alpha < t < \beta\}$ with $0 \leq \alpha < \beta \leq 2\pi$ has measure $\beta - \alpha$.
\end{exercise}

\begin{solution}
The right side defines a finite Borel measure $\nu$, the push-forward of $\lambda$ restricted to $[0, 2\pi]$, since preimages under the continuous $\varphi$ of Borel sets are Borel. For $f = \mathbf{1}_E$,
$\int f\,d\nu = \int_{[0, 2\pi]}f \circ \varphi\,d\lambda$ by definition; by linearity this holds for simple functions, by monotone convergence for measurable $f \geq 0$, and then for $f \in C_c(\RR^2)$. By
Proposition~\ref{measure-theory:proposition-riemann-lebesgue}, $\int_{[0, 2\pi]}f \circ \varphi\,d\lambda = \Lambda f$. So $\nu$ represents $\Lambda$, and $\nu = \sigma$ by the uniqueness in
Theorem~\ref{radon-measures:theorem-riesz-markov}. Since $\varphi$ is injective on $[0, 2\pi)$ and $\varphi(2\pi) = \varphi(0)$, the set of $t \in [0, 2\pi]$ with $\varphi(t)$ in the arc is $(\alpha, \beta)$, of
measure $\beta - \alpha$.
\end{solution}

\begin{exercise}\label{radon-measures:exercise-weak-mass}
Let $\mu_k \rightharpoonup \mu$ be Radon measures on $X$ (locally compact Hausdorff with a countable basis).
\begin{enumerate}
\item Show that $\mu(X) \leq \liminf_k\mu_k(X)$, and that $\mu_k(X) \to \mu(X)$ if $X$ is compact.
\item On $X = \RR$ let $\mu_k = (1 - \frac1k)\delta_0 + \frac1k\delta_k$. Show that $\mu_k \to \delta_0$ narrowly, but $\int x\,d\mu_k = 1$ for all $k$, while
$\int x\,d\delta_0 = 0$: narrow convergence does not control integrals of unbounded functions such as moments.
\end{enumerate}
\end{exercise}

\begin{solution}
(1) The first statement is Proposition~\ref{radon-measures:proposition-weak-convergence-sets}(1) with $U = X$; if $X$ is compact, $1 \in C_c(X)$. (2) For $f \in C_b(\RR)$,
$\int f\,d\mu_k = (1 - \frac1k)f(0) + \frac1kf(k)$, and $|f(k)|/k \leq \|f\|_\infty/k \to 0$, so $\int f\,d\mu_k \to f(0) = \int f\,d\delta_0$. But
$\int x\,d\mu_k = (1 - \frac1k)\cdot0 + \frac1k\cdot k = 1$. The small mass $\frac1k$ escapes to infinity at a rate that keeps the first moment constant; tightness alone does not prevent
this, and one needs a bound such as $\sup_k\int|x|^{1 + \varepsilon}\,d\mu_k < \infty$ (uniform integrability) to pass to the limit in $\int x\,d\mu_k$.
\end{solution}

\begin{exercise}\label{radon-measures:exercise-staircase}
\emph{Length is not continuous.} For $k \geq 1$ let $\gamma_k(x) = (x, \frac1k\sin kx)$ for $x \in [0, 2\pi]$, and let $\mu_k$ be the length measure of $\gamma_k$, the Radon measure on $\RR^2$ with
\[
\int f\,d\mu_k = \int_0^{2\pi}f(\gamma_k(x))\,|\gamma_k^{\prime}(x)|\,dx = \int_0^{2\pi}f\bigl(x, \tfrac1k\sin kx\bigr)\sqrt{1 + \cos^2kx}\,dx .
\]
Show that $\gamma_k \to \gamma_0$, $\gamma_0(x) = (x, 0)$, uniformly, that $\mu_k \rightharpoonup c\mu_0$ with $c = \frac1{2\pi}\int_0^{2\pi}\sqrt{1 + \cos^2t}\,dt$, where $\mu_0$ is the length measure of the
segment, and that $c > 1$. So the curves $\gamma_k$ have length $2\pi c$, which does not converge to the length $2\pi$ of the limit curve.
\end{exercise}

\begin{solution}
$|\gamma_k(x) - \gamma_0(x)| \leq 1/k$. Let $f \in C_c(\RR^2)$, with modulus of uniform continuity $\omega$. Then $|f(\gamma_k(x)) - f(x, 0)| \leq \omega(1/k)$, so $\int f\,d\mu_k$ differs from
$\int_0^{2\pi}f(x, 0)g(kx)\,dx$, where $g(t) = \sqrt{1 + \cos^2t}$, by at most $2\pi\sqrt2\,\omega(1/k)$. The argument of Example~\ref{radon-measures:example-weak-convergence}(5), applied on $[0, 2\pi]$ to the
continuous function $x \mapsto f(x, 0)$ (only its continuity on $[0, 2\pi]$ was used), shows that $\int_0^{2\pi}f(x, 0)g(kx)\,dx \to \bar g\int_0^{2\pi}f(x, 0)\,dx = c\int f\,d\mu_0$. Since $g \geq 1$, with
equality only where $\cos t = 0$, and $g$ is continuous, $c > 1$. The length of $\gamma_k$ is $\int_0^{2\pi}g(kx)\,dx = \frac1k\int_0^{2\pi k}g = 2\pi c$, as the interval consists of $k$ periods. In
general length is only \emph{lower semicontinuous} under uniform convergence of curves: the limit can be shorter, never longer.
\end{solution}

\begin{exercise}\label{radon-measures:exercise-total-variation-point}
Let $L(F) = F_1(0) + 2F_2(0)$ for $F = (F_1, F_2) \in C_c(\RR, \RR^2)$. Find the measure $\mu$ and the field $\nu$ of Theorem~\ref{radon-measures:theorem-riesz-vector}, and check
\eqref{radon-measures:equation-total-variation} directly.
\end{exercise}

\begin{solution}
$L(F) = \langle F(0), (1, 2)\rangle$, so $L(F) = \int\langle F, \nu\rangle\,d\mu$ with $\mu = \sqrt5\,\delta_0$ and $\nu(0) = (1, 2)/\sqrt5$, a unit vector. For an open $V \ni 0$ and $|F| \leq 1$ with
support in $V$, the Cauchy--Schwarz inequality gives $L(F) \leq |(1, 2)| = \sqrt5$, with equality when $F(0) = (1, 2)/\sqrt5$, which is possible for $F = \psi\,(1, 2)/\sqrt5$ with $\psi \prec V$, $\psi(0) = 1$; for
$0 \notin V$ the supremum is $0$. So the total variation is $\sqrt5\,\delta_0$, not $\delta_0 + 2\delta_0 = 3\delta_0$: the components cannot be estimated separately, as $|F| \leq 1$ couples them.
\end{solution}

\begin{exercise}\label{radon-measures:exercise-haar-affine-line}
Let $G$ be the group of Example~\ref{radon-measures:example-ax-plus-b}. Show directly, without the change of variables formula, that $\mu(E) = \int_Ea^{-2}\,da\,db$ is left invariant for rectangles
$E = [a_1, a_2] \times [b_1, b_2]$, and compute $\mu(E(2, 0))$ and $\mu((2, 0)E)$ for $E = [1, 2] \times [0, 1]$.
\end{exercise}

\begin{solution}
Left multiplication by $(a, b)$ maps $(a^{\prime}, b^{\prime})$ to $(aa^{\prime}, ab^{\prime} + b)$, so it maps $[a_1, a_2] \times [b_1, b_2]$ onto $[aa_1, aa_2] \times [ab_1 + b, ab_2 + b]$. By Fubini,
$\mu([a_1, a_2] \times [b_1, b_2]) = (b_2 - b_1)(a_1^{-1} - a_2^{-1})$, and for the image this is $a(b_2 - b_1)\cdot a^{-1}(a_1^{-1} - a_2^{-1})$, the same. For $E = [1, 2] \times [0, 1]$,
$\mu(E) = \frac12$; $(2, 0)E = [2, 4] \times [0, 2]$ has measure $2(\frac12 - \frac14) = \frac12$, while $E(2, 0) = \{(2a, b)\} = [2, 4] \times [0, 1]$ has measure $\frac14 = \Delta(2, 0)\mu(E)$ with
$\Delta(2, 0) = \frac12$.
\end{solution}

\begin{exercise}\label{radon-measures:exercise-haar-finite}
Show that a left Haar measure $\mu$ on $G$ is finite if and only if $G$ is compact, and that $\mu(\{e\}) > 0$ if and only if $G$ is discrete.
\end{exercise}

\begin{solution}
If $G$ is compact, $\mu(G) < \infty$ by definition. If $G$ is not compact, let $V$ be a compact neighbourhood of $e$ with $V = V^{-1}$, and $U$ an open neighbourhood of $e$ with $UU^{-1} \subseteq V$.
Choose inductively $x_1, x_2, \ldots$ with $x_{n+1} \notin x_1V \cup \cdots \cup x_nV$, possible because a finite union of compact sets is not $G$. The sets $x_nU$ are pairwise disjoint: if $x_mu = x_nu^{\prime}$
with $m < n$, then $x_n = x_muu^{\prime-1} \in x_mV$. They have the same measure $\mu(U) > 0$ (Lemma~\ref{radon-measures:lemma-haar-positive}), so $\mu(G) = \infty$. If $G$ is discrete, $\{e\}$ is open and
has positive measure by the same lemma. If $\mu(\{e\}) = c > 0$, then every point has measure $c$, a compact neighbourhood of $e$ has finite measure and so is finite, and $\{e\}$, being the complement in
it of finitely many closed points, is open; so $G$ is discrete.
\end{solution}

\begin{exercise}\label{radon-measures:exercise-haar-quotient-circle}
Show that the Haar probability measure $\sigma$ on $S^1$ satisfies $\int_{S^1}z^k\,d\sigma(z) = 0$ for all integers $k \neq 0$, and deduce that a Borel probability measure $\nu$ on $S^1$ invariant under the
rotation $z \mapsto e^{i\alpha}z$ with $\alpha/2\pi$ irrational equals $\sigma$.
\end{exercise}

\begin{solution}
By invariance under the rotation by $w$, $\int z^k\,d\sigma = \int(wz)^k\,d\sigma = w^k\int z^k\,d\sigma$; choosing $w$ with $w^k \neq 1$ gives $0$. If $\nu$ is invariant under $z \mapsto e^{i\alpha}z$, the
same computation with $w = e^{i\alpha}$, for which $w^k \neq 1$ when $k \neq 0$, gives $\int z^k\,d\nu = 0$ for $k \neq 0$. So $\int p\,d\nu = \int p\,d\sigma$ for every trigonometric polynomial
$p(z) = \sum_{|k| \leq N}a_kz^k$. Trigonometric polynomials are uniformly dense in $C(S^1)$ (by the Stone--Weierstrass theorem, Theorem~\ref{real-analysis:theorem-stone-weierstrass}: they form an algebra closed under conjugation, separating points and containing
constants), so $\int f\,d\nu = \int f\,d\sigma$ for every continuous $f$, and $\nu = \sigma$ by the uniqueness in the Riesz--Markov theorem.
\end{solution}
